% !TeX program = pdflatex
% Apery-like recurrences beyond the rigid case:
% four-term rows, K3 surfaces, and new periods.
% Self-contained; builds with pdflatex; bibliography inline.
\documentclass[11pt,reqno]{amsart}

\usepackage[T1]{fontenc}
\usepackage[utf8]{inputenc}
\usepackage{amsmath,amssymb,amsthm}
\usepackage{mathtools}
\usepackage{array}
\usepackage{booktabs}
\usepackage{enumitem}
\usepackage[margin=0.85in]{geometry}
\usepackage{microtype}
\usepackage[hidelinks]{hyperref}

\theoremstyle{plain}
\newtheorem{theorem}{Theorem}[section]
\newtheorem{proposition}[theorem]{Proposition}
\newtheorem{lemma}[theorem]{Lemma}
\newtheorem{corollary}[theorem]{Corollary}
\newtheorem{conjecture}[theorem]{Conjecture}
\theoremstyle{definition}
\newtheorem{definition}[theorem]{Definition}
\theoremstyle{remark}
\newtheorem{remark}[theorem]{Remark}

\setlength{\emergencystretch}{3em}
\setlength{\abovedisplayskip}{4pt plus 2pt minus 2pt}
\setlength{\belowdisplayskip}{4pt plus 2pt minus 2pt}
\setlength{\abovedisplayshortskip}{2pt plus 2pt}
\setlength{\belowdisplayshortskip}{2pt plus 2pt}
\setlength{\textfloatsep}{8pt plus 2pt minus 2pt}
\setlength{\floatsep}{8pt plus 2pt minus 2pt}
\setlength{\intextsep}{8pt plus 2pt minus 2pt}
\makeatletter
\def\thm@space@setup{\thm@preskip=4pt \thm@postskip=4pt}
\makeatother
\renewcommand{\topfraction}{0.92}
\renewcommand{\bottomfraction}{0.8}
\renewcommand{\textfraction}{0.05}
\renewcommand{\floatpagefraction}{0.85}

\theoremstyle{plain}
\newtheorem{namedthm}{Theorem}
\renewcommand{\thenamedthm}{\Alph{namedthm}}

\newcommand{\Z}{\mathbb{Z}}
\newcommand{\Q}{\mathbb{Q}}
\newcommand{\R}{\mathbb{R}}
\newcommand{\C}{\mathbb{C}}
\newcommand{\PP}{\mathbb{P}}
\DeclareMathOperator{\PSL}{PSL}
\DeclareMathOperator{\Gal}{Gal}
\DeclareMathOperator{\Res}{Res}
\DeclareMathOperator{\re}{Re}
\DeclareMathOperator{\im}{Im}
\newcommand{\Jm}{\mathcal J}
\newcommand{\Rc}{\mathcal R}
\newcommand{\Sc}{\mathcal S}
\newcommand{\Vc}{\mathcal V}
\newcommand{\Ic}{\mathcal I}

% ---------------------------------------------------------------------------
% Evidence-class markers.  These are load-bearing and are never upgraded.
% ---------------------------------------------------------------------------
\newcommand{\Proved}{\textup{\textsf{[PROVED]}}}
\newcommand{\Verified}[1]{\textup{\textsf{[VERIFIED: #1]}}}
\newcommand{\Numerical}[1]{\textup{\textsf{[NUMERICAL: #1]}}}
\newcommand{\Conj}{\textup{\textsf{[CONJECTURE]}}}
\newcommand{\Open}{\textup{\textsf{[OPEN]}}}

\renewcommand{\arraystretch}{0.96}
\begin{document}

\title[Beyond the rigid case: four-term rows]
      {Ap\'ery-like recurrences beyond the rigid case:\\
       four-term rows, $K3$ surfaces, and new periods}
\author{Claude (Fable), with River}
\date{\today}

\begin{abstract}
Integral three-term Ap\'ery-like recurrences $(n+1)^2u_{n+1}=P(n)u_n-Q(n)u_{n-1}$
have a rank-two Picard--Fuchs operator with four singular points: a rigid
situation, classified in the companion paper under the hypothesis that the local
system comes from an elliptic surface. We take the first step out of rigidity,
to the four-term row
$(n+1)^2u_{n+1}=P(n)u_n-Q(n)u_{n-1}+R(n)u_{n-2}$, whose operator has rank two
with \emph{five} singular points and two accessory parameters. We give the
exponent dictionary and the normalisation classes, both in the equal-exponent
and in the mixed-exponent case, the Riemann--Hurwitz bound
$\deg\Jm\le6c+4e_3+3e_2-12\le18$, and the results of complete or box-complete
scans. Four things happen at once. \emph{Near-total rigidity persists}: of
$2178$ integral rows produced, $98.2\%$ have an apparent singularity and are
gauge or cusp-move images of three-term rows, and the only one-parameter
families found are of that kind. \emph{A finite sporadic list appears}: $30$
primitive rows with five genuine singular points, $28$ of them with denominator
exponent $k=2$. \emph{$K3$ geometry enters exactly here}: one of the $30$ is the
Picard--Fuchs system of an elliptic $K3$ with five singular fibres
$I_4I_6I_3I_3IV^*$, which is impossible with four singular points. \emph{New
periods appear}: in the mixed-exponent classes there are genuine five-point rows
with irrational conjugate singularities whose Ap\'ery limits are $\tfrac14G$,
$\tfrac12G-\tfrac3{16}\zeta(2)$, $\tfrac38\zeta(2)-\tfrac12G$ and
$L(2,\chi_{-3})$-combinations, whose gauge partners all give
$\Gamma(\tfrac14)^4/(64\pi)$, and, one term further, a six-point five-term row
whose Ap\'ery limit is the weight-$3$ cusp-form value $\tfrac{\sqrt2}3L(g,2)$.
What does \emph{not} appear is new irrationality-relevant arithmetic: no genuine
five-point row has positive score, no unequal-exponent class can contain a
decaying row, and the $2$-adic content of the new Catalan rows is the ledger's
own.
\end{abstract}

\maketitle

% ===========================================================================
\section{Introduction}\label{sec:intro}
% ===========================================================================

\subsection{The rigid case, and one step out}
Let
\begin{equation}\label{eq:three}
(n+1)^2u_{n+1}=P(n)u_n-Q(n)u_{n-1},\qquad u_0=1,\ u_{-1}=0,
\end{equation}
with $P,Q$ integral of degree $\le2$. The generating function $y=\sum u_nt^n$
satisfies a rank-two Fuchsian equation with the four singular points
$0,t_1,t_2,\infty$, and such an equation is \emph{rigid}: no accessory parameter,
local exponents determining it, and --- when the local system comes from an
elliptic surface --- the surface forced onto Beauville's and Herfurtner's finite
lists \cite{Beauville1982,Herfurtner1991}. The companion paper \cite{CLASS}
carries this out: its Theorem~A is a conditional classification (all four fibres
semistable, $\deg\Jm\le12$ by Riemann--Hurwitz and $12\mid\deg\Jm$ by a
twist-parity lemma, hence $\deg\Jm=12$ and a \emph{rational} elliptic surface,
with the integral rows exactly Zagier's seven \cite{Zagier2009}), its Theorem~B
collects the unconditional necessary conditions, and the missing converse is
Zagier's conjecture that every integral row has a modular parametrisation ---
Bombieri--Dwork in rank two with four singular points.

This paper asks what survives one step out. The object is the integral
\emph{four-term} row
\begin{equation}\label{eq:four}
(n+1)^2u_{n+1}=P(n)u_n-Q(n)u_{n-1}+R(n)u_{n-2},\qquad u_0=1,\ u_{-1}=u_{-2}=0,
\end{equation}
$P=an^2+bn+c$, $Q=dn^2+en+f$, $R=gn^2+hn+j$ integral with $g\ne0$. Its operator
is rank two with \emph{five} singular points $0,t_1,t_2,t_3,\infty$ and
therefore has $5-3=2$ accessory parameters --- it is the first non-rigid case.
Geometrically the change is just as sharp: elliptic surfaces with four singular
fibres form finitely many configurations, while those with five come in
one-parameter families (Persson \cite{Persson1990}), so the cross-ratio
finiteness mechanism of the rigid case disappears. Longer recurrences and their
accessory parameters are of course not new --- they occur throughout the
Ap\'ery-like literature \cite{AlmkvistZudilin2006,Cooper2012,MovasatiReiter} and
the accessory-parameter problem for such operators is Dwork's
\cite{Beukers2002} --- but the systematic question of which \emph{integral} rows
exist, and what they compute, has not been asked one step beyond rank two with
four points.

\subsection{The answers, in words}
Four things happen, and one does not.

\begin{enumerate}[label=\textup{(\arabic*)},leftmargin=2.4em]
\item \emph{Near-total rigidity persists.} Of the $2178$ distinct integral rows
the census produces, $2138$ ($98.2\%$) have an \emph{apparent} singularity --- a
point where the exponent difference is a positive integer and the Frobenius
obstruction vanishes, or a half-integral one where the monodromy is $-I$. Their
\emph{projective} local systems have only four singular points, so they are gauge
or cusp-move images of three-term rows. Every one-parameter family the scan
produced is of this kind, and the largest is written down in closed form
(Proposition~\ref{prop:F4}): the signed binomial transform
$v_n=\sum_m\binom nm(-\nu)^{n-m}u_m$ of a Zagier-class row at a
non-characteristic $\nu$.

\item \emph{A finite sporadic list appears.} After deduplication, $30$ primitive
rows with five genuine singular points remain (Table~\ref{tab:thirty}); $28$ of
them have sharp denominator exponent $k=2$, the four-term echo of the free
integration phenomenon, and the other two are Pad\'e-type with $k=1$. No
one-parameter integral family of genuine five-point rows was found.

\item \emph{$K3$ geometry enters exactly here.} With four singular points
$\deg\Jm\le12$ always, so no elliptic $K3$ can be the Picard--Fuchs system of a
three-term row. With five the bound is $18$, and one row of the sporadic list
attains $\deg\Jm=16$ and is an elliptic $K3$ with fibres $I_4I_6I_3I_3IV^*$,
Picard number $20$ and transcendental motive the weight-$3$ CM newform
\texttt{32.3.d.a}. That row is the subject of the companion note \cite{K3PAPER};
we state it (\S\ref{sec:k3}) because it is what the scan found, without repeating
its proofs.

\item \emph{New periods appear.} The mixed-exponent classes --- exponent
differences $(0,\tfrac12,\tfrac12)$ at the three finite points, which the first
census never scanned --- contain genuine five-point rows with \emph{irrational
conjugate} singular points whose Ap\'ery limits are $\tfrac14G$,
$\tfrac12G-\tfrac3{16}\zeta(2)$, $\tfrac38\zeta(2)-\tfrac12G$ and three
$L(2,\chi_{-3})$-combinations, with the whole period matrix classical; the same
class in the other half-integral gauge gives $\Gamma(\tfrac14)^4/(64\pi)$ on
three different rows. One term further still, a six-point five-term row has the
weight-$3$ cusp-form value $\tfrac{\sqrt2}3L(g,2)$ as an honest Ap\'ery limit ---
the $L$-value that the $K3$ row itself, having no archimedean limit, can produce
only as a fold constant.
\end{enumerate}

What does not happen is arithmetic gain. No genuine five-singular-point row in
any scanned class has positive irrationality score $\log(1/|\lambda_2|)-k$; the
only positive scores belong to disguised three-term rows with $k=1$, as in the
rigid case. Structurally, a row with $|\lambda_2|<1$ must have an
\emph{irreducible} characteristic cubic (Theorem~\ref{thm:D1}), so no
unequal-exponent class can contain one, and a \emph{real} such $\lambda_2$ forces
the cubic to be totally real. The new Catalan-period rows realise the geometric
evasion (E1) of the obstruction analysis \cite{CATOBS} --- irrational conjugate
singularities, $+0.693$ nats of archimedean score over Zagier $\mathbf E$ ---
but their $2$-adic slope set is $\{2\}$ and their limits lie in the ledger's own
$\zeta_2(2)$ and $L_2(2,\chi_{12})$ classes; every one is dominated, as a lattice
engine, by the row it would replace.

\subsection{Evidence classes}
Statements carry one of the markers \Proved{} (an exact argument),
\Verified{range} (an exact finite computation over the stated range: exact
rational or $p$-adic arithmetic, or a high-precision computation whose agreement
is stated), \Numerical{precision} (a floating-point measurement), \Conj{} and
\Open, and these are never upgraded. All scans reported are $a$-block ordered, so
every reported box is a complete scan of the stated prefix. Sources, scripts and
logs are named in \S\ref{sec:repro}.

\subsection{Main statements}
For orientation we record the three named theorems that organise the paper; each
is proved or verified in the section indicated.

\begin{namedthm}[the dictionary and its classes]\label{thmA}
\Proved{} \textup{(\S\ref{sec:dict})} The operator of \eqref{eq:four} is
$t^2\Rc D^2+t\Sc D+t\Vc$ with $\Rc,\Sc,\Vc$ as in \eqref{eq:RSV}, and its
exponents are $(0,0)$ at $t=0$, $(0,\rho_i)$ at the roots $t_i$ of $\Rc$ with
$\rho_i=-T(t_i)/(t_i\Rc'(t_i))$, $T=\Sc-t\Rc'$, and $(2-s_1,2-s_2)$ at $\infty$
with $s_1,s_2$ the roots of $R$; Fuchs reads $\sum\rho_i=s_1+s_2-1$. All three
$\rho_i$ coincide if and only if three \emph{linear} conditions on
$(b,e,h)$ hold, and that is automatic when the characteristic cubic is
irreducible; the unequal case with a quadratic Galois orbit has the
normalisation forms of Theorem~\ref{thm:D3}. If the local system is that of an
elliptic surface then all exponent differences lie in
$\{0,\tfrac12,\tfrac13,\tfrac23\}\bmod1$ and
$\deg\Jm\le6c+4e_3+3e_2-12\le18$.
\end{namedthm}

\begin{namedthm}[the census]\label{thmB}
\Verified{boxes as stated in \S\ref{sec:scans}} Over $27$ Kodaira-admissible
equal-exponent classes the scan produces $2178$ distinct integral rows: $4$ fail
exact integrality at $n=120$, $2138$ ($98.2\%$) have an apparent singularity and
are disguised three-term rows, $6$ are rescalings, and $30$ are primitive rows
with five genuine singular points, exactly one of which passes the $\Jm$-map test
and is an elliptic $K3$. The two productive mixed-exponent classes, swept to
completion, contain exactly nine primitive genuine five-point rows each, in
twist-partner bijection.
\end{namedthm}

\begin{namedthm}[no arithmetic gain]\label{thmC}
\Proved{} and \Verified{boxes W, PW and the extended window} A row with
$|\lambda_2|<1$ has an irreducible characteristic cubic, and a real such
$\lambda_2$ forces it to be totally real; consequently no unequal-exponent class
contains a row with $|\lambda_2|<1$. The complete positive-score window over all
$27$ equal-exponent classes contains exactly two integral rows, both with $k=1$
and both carrying an $I_0^*$ at $\infty$, hence disguised: \emph{no genuine
five-singular-point four-term row has a positive score}.
\end{namedthm}

% ===========================================================================
\section{The dictionary and the normalisation classes}\label{sec:dict}
% ===========================================================================

\subsection{The operator}
Write $\theta=t\,d/dt$, $D=d/dt$.

\begin{lemma}[\Proved]\label{lem:op}
$y=\sum_{n\ge0}u_nt^n$ satisfies $Ly=0$ with
$L=\theta^2-t\,P(\theta)+t^2Q(\theta+1)-t^3R(\theta+2)$.
\end{lemma}

\begin{proof}
$[t^{n+1}]\theta^2y=(n+1)^2u_{n+1}$, $[t^{n+1}]tP(\theta)y=P(n)u_n$,
$[t^{n+1}]t^2Q(\theta+1)y=Q(n)u_{n-1}$ and
$[t^{n+1}]t^3R(\theta+2)y=R(n)u_{n-2}$.
\end{proof}

Expanding $\theta=tD$, $\theta^2=t^2D^2+tD$ gives
$L=t^2\Rc(t)D^2+t\,\Sc(t)D+t\,\Vc(t)$ with
\begin{equation}\label{eq:RSV}
\begin{aligned}
\Rc(t)&=1-at+dt^2-gt^3,\\
\Sc(t)&=1-(a+b)t+(3d+e)t^2-(5g+h)t^3,\\
\Vc(t)&=-c+(d+e+f)t-(4g+2h+j)t^2 .
\end{aligned}
\end{equation}
The singular points are $t=0$, the three roots $t_1,t_2,t_3$ of $\Rc$, and
$t=\infty$. Since $t^3(\lambda^3-a\lambda^2+d\lambda-g)=\Rc(1/\lambda)$ with
$\lambda=1/t$, the \emph{characteristic roots} $\lambda_i=1/t_i$ are the roots of
\begin{equation}\label{eq:chi}
\chi(\lambda)=\lambda^3-a\lambda^2+d\lambda-g ,
\end{equation}
and five distinct singular points requires $g\ne0$ and
$\Delta=18adg-4a^3g+a^2d^2-4d^3-27g^2\ne0$.

\begin{theorem}[exponent dictionary; \Proved]\label{thm:F1}
The exponents of $L$ are $(0,0)$ at $t=0$ \textup{(}a MUM point, so a logarithm
is forced\textup{)}; $(0,\rho_i)$ at a simple root $t_i$ of $\Rc$, with
\[
T:=\Sc-t\Rc'=1-bt+(d+e)t^2-(2g+h)t^3,\qquad
\rho_i=1-\frac{\Sc(t_i)}{t_i\Rc'(t_i)}=-\frac{T(t_i)}{t_i\Rc'(t_i)} ;
\]
and $(2-s_1,2-s_2)$ at $t=\infty$, where $s_1,s_2$ are the roots of $R$ and
$\delta_\infty:=s_2-s_1$. Moreover $\sum_i\rho_i=s_1+s_2-1=-h/g-1$.
\end{theorem}

\begin{proof}
The first two are the indicial equations of $t^2\Rc D^2+t\Sc D+t\Vc$ at $t=0$ and
at a simple zero of $\Rc$. At $\infty$, $Lt^{\mu}$ has leading term
$-R(\mu+2)t^{\mu+3}$, so $\mu+2\in\{s_1,s_2\}$; with the convention
(exponent)$\,=-\mu$ the exponents are $2-s_i$. For the last claim,
$\sum_i\rho_i=-\sum_i\Res_{t_i}\bigl(T/(t\Rc)\bigr)$; the rational function
$T/(t\Rc)$ has residue $T(0)/\Rc(0)=1$ at $t=0$ and behaves like
$\bigl((2g+h)/g\bigr)t^{-1}$ at $\infty$, so the sum of its finite residues is
$(2g+h)/g$ and $\sum\rho_i=1-(2g+h)/g$. The sum of all ten exponents is then
$3=(5-2)\cdot1$, the Fuchs relation for rank two with five singular points.
\end{proof}

The identities of Lemma~\ref{lem:op} and Theorem~\ref{thm:F1} are also checked
symbolically \Verified{$n\le12$, script \texttt{01\_exponents.py}}.

\begin{remark}[the accessory parameters and the M\"obius invariants]
The scaling $t\mapsto t/\mu$, i.e.\ $u_n\mapsto\mu^nu_n$, acts by
$(a,b,c)\mapsto\mu(a,b,c)$, $(d,e,f)\mapsto\mu^2(d,e,f)$,
$(g,h,j)\mapsto\mu^3(g,h,j)$, so the labelled configuration
$(0,\infty;\{t_1,t_2,t_3\})$ has exactly two invariants,
\[
\Ic_1=\frac{a^3}{g}=\frac{(\lambda_1+\lambda_2+\lambda_3)^3}{\lambda_1\lambda_2\lambda_3},
\qquad
\Ic_2=\frac{a^2}{d}=\frac{(\lambda_1+\lambda_2+\lambda_3)^2}{\sum_{i<j}\lambda_i\lambda_j},
\]
the analogue of the cross-ratio invariant of \cite{CLASS}. Both are rational for
an integral row --- and, unlike the four-point case, this is no longer
restrictive: the realisable locus for a five-fibre configuration is a
\emph{curve} in the $(\Ic_1,\Ic_2)$-plane, not a point, those configurations
coming in one-parameter families \cite{Persson1990}. That is the precise sense in
which the four-term case is non-rigid, and why the $\Jm$-map test has to do all
the work later.
\end{remark}

\subsection{Equal exponents}

\begin{theorem}[normalisation classes; \Proved]\label{thm:F2}
Assume $t_1,t_2,t_3$ distinct. Then $\rho_1=\rho_2=\rho_3=\rho$ if and only if
\[
b=(1-\rho)a,\qquad e=-2\rho\,d,\qquad h=-(1+3\rho)\,g .
\]
\end{theorem}

\begin{proof}
$\rho_i=\rho$ for all $i$ says $\Rc\mid T+\rho\,t\Rc'$. Both are cubics with
constant term $1$, so $T+\rho t\Rc'=\Rc$, i.e.
$1-(b+\rho a)t+(d+e+2\rho d)t^2-(2g+h+3\rho g)t^3=1-at+dt^2-gt^3$.
\end{proof}

\begin{corollary}[Galois; \Proved]\label{cor:galois}
The $\rho_i$ are values at $t_i$ of a fixed rational function with rational
coefficients, hence are permuted by $\Gal(\overline\Q/\Q)$ exactly as the $t_i$
are. Kodaira admissibility \textup{(}Corollary~\ref{cor:kodaira}\textup{)} forces
$\rho_i\in\Q$; therefore $\rho_i$ is constant on each Galois orbit of
$\{t_1,t_2,t_3\}$, and in particular if $\chi$ is irreducible then
$\rho_1=\rho_2=\rho_3$.
\end{corollary}

\begin{corollary}[Kodaira admissibility; \Proved{} given Kodaira's list]\label{cor:kodaira}
If the local system of $L$ is, up to a rational gauge factor and a quadratic
twist, $R^1\pi_*\Q$ of an elliptic surface over $\PP^1_t$, then
$\rho_1,\rho_2,\rho_3,\delta_\infty\in\{0,\tfrac12,\tfrac13,\tfrac23\}\pmod1$.
\end{corollary}

\begin{proof}
Neither a rational gauge nor a quadratic twist changes an exponent
\emph{difference} mod $1$, and Kodaira's monodromy eigenvalue ratios are $1$
($I_n,I_n^*$), $-1$ ($III,III^*$), $e^{\pm2\pi i/3}$ ($IV,IV^*$),
$e^{\mp2\pi i/3}$ ($II,II^*$) \cite{Kodaira1963,MirandaBook}. This is the
argument of \cite[Theorem~B(ii)]{CLASS} verbatim.
\end{proof}

\subsection{The class parametrisation}
Fix $\rho$ and write $R(n)=C(Mn-j_1)(Mn-j_2)$, so $s_i=j_i/M$ and
$g=CM^2$, $h=-CM(j_1+j_2)$, $j=Cj_1j_2$, with $j_1+j_2=(1+3\rho)M$ by
Theorems~\ref{thm:F1} and \ref{thm:F2}. A \emph{class} is the tuple
$(\rho;M,j_1,j_2)$, and inside it the scan runs over
\begin{equation}\label{eq:class}
P(n)=a\bigl(n^2+(1-\rho)n\bigr)+c,\qquad
Q(n)=d\bigl(n^2-2\rho n\bigr)+f,\qquad
R(n)=C(Mn-j_1)(Mn-j_2)
\end{equation}
with the five integer parameters $a,c,d,f,C$ (four after $t\mapsto t/\mu$). The
two accessory parameters are visible as $c=u_1=P(0)$ and $f=Q(0)$. There are
$156$ Kodaira-admissible classes with $\rho,\delta_\infty$ of denominator $\le6$,
$\rho\in[-1,2]$, $\delta_\infty\in[0,4]$ and $M\le12$
\Verified{\texttt{04\_classes.py}}.

\subsection{The semistable classes}
``All five exponent differences $\equiv0$'' means $\rho\in\Z$ and
$\delta_\infty\in\Z$. Since $s_1+s_2=1+3\rho$, and $\delta_\infty=0$ forces
$s_1=s_2=(1+3\rho)/2$, the exponents at $\infty$ are integral iff $\rho$ is odd.
Three presentations matter:
\[
\begin{array}{llll}
(\mathbf S_0)&P=a(n^2+n)+c,&Q=dn^2+f,&R=C(2n-1)^2,\\
(\mathbf S_1)&P=an^2+c,&Q=d(n^2-2n)+f,&R=C(n-2)^2,\\
(\mathbf S_{-1})&P=a(n^2+2n)+c,&Q=d(n^2+2n)+f,&R=C(n+1)^2 .
\end{array}
\]
$(\mathbf S_0)$ has exponents $(\tfrac32,\tfrac32)$ at $\infty$: four $I_n$
fibres and one $I_m^*$. $(\mathbf S_{\pm1})$ has integral exponents at $\infty$
and five $I_n$ fibres; the two are gauge partners, the gauge
$y\mapsto\prod_i(1-\lambda_it)^{-\rho}y$ sending $\rho\mapsto-\rho$ and
$s\mapsto s-3\rho$. The distinction between $\mathbf S_0$ and $\mathbf S_{\pm1}$
is not a matter of gauge: a quadratic character twist flips the starredness of
the fibres at an \emph{even} number of points, so the parity of the number of
$I^*$ fibres is an invariant of the projective local system --- $1$ for
$\mathbf S_0$, $0$ for $\mathbf S_{\pm1}$. This is the four-term form of the
twist-parity lemma of \cite{CLASS}.

$(\mathbf S_0)$ is the exact analogue of Zagier's shape $(an^2+an+b,\;cn^2)$:
the linear conditions are $b=a$, $e=0$, $h=-g$, $j=g/4$, and both accessory
parameters are free --- unlike the three-term case, $f$ is \emph{not} forced to
vanish, the fibre at $\infty$ being governed by $R$, not $Q$. It carries integral
rows; $(\mathbf S_{\pm1})$ does not (Proposition~\ref{prop:F5}).

\subsection{Apparent singularities}\label{sec:apparent}
At $t_i$ the exponents are $(0,\rho_i)$. If $\rho_i\in\Z\setminus\{0\}$ the two
exponents differ by an integer and the second solution either carries a logarithm
(a genuine $I_n$ fibre) or does not, in which case the local monodromy is
trivial: an \emph{apparent} singularity. The same happens at $\infty$ when
$\delta_\infty\in\Z_{>0}$. If $\rho_i=0$ (or $\delta_\infty=0$) the exponents
coincide and a logarithm is unavoidable, so the point is always a genuine $I_n$.

The obstruction is a single Frobenius coefficient: writing the equation at the
point as $Ay''+By'+Cy=0$ and running the recursion from the smaller exponent, the
coefficient at $m=|\rho_i|$ (resp.\ $m=\delta_\infty$) has vanishing indicial
factor, and the point is apparent iff the accumulated right-hand side vanishes
\Verified{$100$ digits, relative threshold $10^{-40}$}. Two cases must be
distinguished: with integral exponents the local monodromy is the identity and
the point is not a singular point of the local system at all; with half-integral
ones --- which happens at $\infty$ when $\delta_\infty\in\Z$ but $s_i\notin\Z$ ---
it is $-I$, trivial in $\PSL_2$, i.e.\ an $I_0^*$ fibre. Both are invisible to
the \emph{projective} local system, the one the $\Jm$-map test and the Kodaira
dictionary see, so in both cases the row counts as a four-point object --- the
$I_0^*$ mechanism that produced the one-parameter family classes in the rigid
case \cite{CLASS}.

\begin{proposition}[the class $\mathbf S_{-1}$; \Proved]\label{prop:F5}
In the class $(\rho;\delta_\infty)=(-1;0)$ the Frobenius obstruction at $t_i$ is
$S_i=t_i\bigl[(a-c)+(f-d)t_i\bigr]$. Hence at most one $t_i$ is apparent unless
$a=c$ and $d=f$, and in that case all three are, which happens exactly when
$P(n)=a(n+1)^2$, $Q(n)=d(n+1)^2$, $R(n)=C(n+1)^2$, i.e.\ exactly when the row
degenerates to the constant-coefficient recurrence
$u_{n+1}=au_n-du_{n-1}+Cu_{n-2}$ and $\sum u_nt^n$ is a rational function.
\end{proposition}

\begin{proof}
Substitute $b=2a$, $e=2d$, $h=2g$, $j=g$ into $S_i=2A_2-B_1+C_0$ with
$A=t^2\Rc$, $B=t\Sc$, $C=t\Vc$ expanded at $t_i$, and use $\Rc(t_i)=0$ and
$\Sc(t_i)=2t_i\Rc'(t_i)$ (which is $\rho_i=-1$); this gives
$S_i=t_i[2\Rc'+t_i\Rc''-\Sc'+\Vc](t_i)=t_i[(a-c)+(f-d)t_i]$ identically in
$t_i$.
\end{proof}

\subsection{Mixed exponents}\label{sec:mixdict}
Corollary~\ref{cor:galois} confines unequal exponents to reducible cubics, and
the general such shape has a rational root and a quadratic Galois orbit. Let
$W(\lambda)=\lambda^3+\beta_1\lambda^2+\beta_2\lambda+\beta_3$ be the reversal of
$T$, i.e.\ $T(t)=1+\beta_1t+\beta_2t^2+\beta_3t^3$.

\begin{lemma}[\Proved]\label{lem:D30}
$W(\lambda)=\chi(\lambda)\bigl[1+\sum_i\rho_i\lambda_i/(\lambda-\lambda_i)\bigr]
 =\chi(\lambda)+\sum_i\rho_i\lambda_i\prod_{j\ne i}(\lambda-\lambda_j)$,
and $b=-\beta_1$, $d+e=\beta_2$, $2g+h=-\beta_3$.
\end{lemma}

\begin{proof}
With $\Rc=\prod_i(1-\lambda_it)$ one has
$t_i\Rc'(t_i)=-(\lambda_i-\lambda_j)(\lambda_i-\lambda_k)/\lambda_i^2$, so
$\rho_i=-T(t_i)/(t_i\Rc'(t_i))$ reads $W(\lambda_i)=\rho_i\lambda_i\chi'(\lambda_i)$.
$W$ is the unique \emph{monic} cubic with those three values; Lagrange
interpolation with $L_i(\lambda)=\chi(\lambda)/((\lambda-\lambda_i)\chi'(\lambda_i))$
gives the display, the leading coefficient being $1$ because the correction term
has degree $\le2$.
\end{proof}

Setting all $\rho_i=\rho$ recovers $W=\rho\lambda\chi'+(1-3\rho)\chi$, i.e.\
Theorem~\ref{thm:F2}.

\begin{theorem}[mixed normalisation forms; \Proved]\label{thm:D3}
Let $\chi=(\lambda-r)(\lambda^2-\sigma\lambda+\pi)$ with $r,\sigma,\pi\in\Z$,
and put $\rho_r$ on the rational root and $\rho_p$ on both roots of the quadratic
factor. Then $a=r+\sigma$, $d=\sigma r+\pi$, $g=r\pi$ and
\[
b=(1-\rho_r)r+(1-\rho_p)\sigma,\qquad
e=-\rho_p(2\pi+r\sigma)-\rho_r r\sigma,\qquad
h=-(1+2\rho_p+\rho_r)\,g,
\]
and, writing $R(n)=C(Mn-j_1)(Mn-j_2)$ as in \eqref{eq:class},
$g=CM^2$, $j=Cj_1j_2$ and $(j_1+j_2)/M=s_1+s_2=1+2\rho_p+\rho_r$. Fuchs is
automatic: $\sum\rho_i=2\rho_p+\rho_r=s_1+s_2-1$.
\end{theorem}

\begin{proof}
Expand Lemma~\ref{lem:D30} with $\rho_1=\rho_r$ at $\lambda_1=r$ and
$\rho_2=\rho_3=\rho_p$, using
$\lambda_2\prod_{j\ne2}(\lambda-\lambda_j)+\lambda_3\prod_{j\ne3}(\lambda-\lambda_j)
=(\lambda-r)(\sigma\lambda-2\pi)$, so that
$W=\chi+\rho_p(\lambda-r)(\sigma\lambda-2\pi)+\rho_r r(\lambda^2-\sigma\lambda+\pi)$,
and read off the coefficients.
\end{proof}

Theorem~\ref{thm:D3} was checked numerically against the dictionary of
Theorem~\ref{thm:F1} for all $15$ shapes tried, six random parameter points each,
$0$ failures \Verified{\texttt{01\_mixcheck.py}}. A mixed class is a tuple
$(\rho_p,\rho_r;M,j_1,j_2)$; twenty are Kodaira-admissible, non-Casoratian and
genuinely mixed:

The twenty are: $(\rho_p,\rho_r)=(\pm\tfrac12,0)$, fibres $I_n,III,III$, with
$(M,j_1,j_2)$ equal to $(4,3,5)$, $(6,5,7)$, $(3,2,4)$, $(2,1,3)$ resp.\
$(1,0,0)$, $(4,-1,1)$, $(6,-1,1)$, $(3,-1,1)$; $(0,\pm\tfrac12)$, fibres
$III,I_n,I_n$, with $(4,3,3)$, $(12,7,11)$, $(12,5,13)$, $(2,0,3)$ resp.\
$(4,1,1)$, $(2,0,1)$, $(12,1,5)$, $(12,-1,7)$; $(\pm\tfrac12,\mp\tfrac12)$,
three $III$ fibres in two gauges, with $(4,3,3)$ and $(4,1,1)$; and the two
controls $(0,\pm1)$, where apparent points are expected.


Because $\rho_p=\pm\tfrac12\notin\Z$ and $\rho_r=0$ forces a logarithm, the
classes in the first two rows admit \emph{no} apparent singularity at a finite
point --- only $\infty$ can be apparent, and only when $\delta_\infty\in\Z$ ---
so they carry genuine five-point rows automatically. Structurally they are the
quadratic twists, ramified at exactly two of the three finite points, of the
equal-exponent classes: the twist turns two $I_n$ fibres into $I_n^*$ and leaves
the \emph{projective} local system unchanged. They are not a new geometry but a
new \emph{integral form} of the same geometry, and, as for the $\pm\tfrac12$
twist pairs of \S\ref{sec:thirty} note~3, their Ap\'ery limits differ from those
of the untwisted rows.

\subsection{The Riemann--Hurwitz window}\label{sec:RH}

\begin{theorem}[\Proved]\label{thm:F3}
Let $\pi:\mathcal E\to\PP^1_t$ be a minimal elliptic surface with section and
non-constant $\Jm$ whose singular fibres lie over the five singular points of
$L$, and let $c,e_2,e_3$ count the points whose projective local monodromy is
parabolic, of order $2$, of order $3$ \textup{(}so $c+e_2+e_3\le5$ and $c\ge1$
because $t=0$ is MUM\textup{)}. Then
$\mu=\deg\Jm\le6c+4e_3+3e_2-12\le18$.
\end{theorem}

\begin{proof}
Riemann--Hurwitz for $\Jm:\PP^1\to\PP^1$ together with Kodaira's dictionary gives
$2\mu-2\ge(\mu-c)+\tfrac23(\mu-e_3)+\tfrac12(\mu-e_2)$; maximising subject to
$c+e_2+e_3\le5$ gives $\mu\le6\cdot5-12=18$, attained only in the all-parabolic
case. The argument is that of the four-point bound in \cite{CLASS}, which never
used the number of singular points.
\end{proof}

By Gauss--Bonnet a genus-zero group $\Gamma\le\PSL_2(\Z)$ with exactly five
special points ($c$ cusps, $e_2$ elliptic of order $2$, $e_3$ of order $3$,
$c+e_2+e_3=5$, $c\ge1$) has index
$\mu=12\bigl(-1+\tfrac c2+\tfrac{e_2}4+\tfrac{e_3}3\bigr)=6c+4e_3+3e_2-12$,
which is exactly the bound of Theorem~\ref{thm:F3}: the two agree because
equality says precisely that $\Jm$ is a Belyi map. The fifteen possible
signatures give indices
$(5,0,0)\!:\!18$, $(4,1,0)\!:\!15$, $(4,0,1)\!:\!16$, $(3,2,0)\!:\!12$,
$(3,1,1)\!:\!13$, $(3,0,2)\!:\!14$, $(2,3,0)\!:\!9$, $(2,2,1)\!:\!10$,
$(2,1,2)\!:\!11$, $(2,0,3)\!:\!12$, $(1,4,0)\!:\!6$, $(1,3,1)\!:\!7$,
$(1,2,2)\!:\!8$, $(1,1,3)\!:\!9$, $(1,0,4)\!:\!10$,
so there are \emph{finitely many host groups} --- the five-point analogue of the
$28$ genus-zero groups with four special points in the rigid case. Their
enumeration into conjugacy classes is not carried out here \Open. Accordingly
the $\Jm$-map test is run with $\mathrm{DMAX}=\mathrm{HMAX}=18$, $130$ series
terms and $760$ values of $\gamma$ per $h$.

% ===========================================================================
\section{The scans}\label{sec:scans}
% ===========================================================================

\subsection{Design}
The scanner runs one class $(\rho;M,j_1,j_2)$ at a time over the five integer
parameters $a,c,d,f,C$ of \eqref{eq:class}. Integrality is tested
division-free on $U_n=u_n(n!)^2$,
\[
U_{n+1}=P(n)U_n-n^2Q(n)U_{n-1}+n^2(n-1)^2R(n)U_{n-2},\qquad
u_n\in\Z\iff v_p(U_n)\ge2v_p(n!)\ \forall p,
\]
prime by prime modulo $p^K$ with $K=2v_p(N!)+2$, $N=24$. Three cheap layers
precede the exact test: analytic congruences ($U_2,\dots,U_4$ --- and, in the
deep scans, $U_5$ --- are \emph{linear in $C$}, so $4\mid U_2$, $36\mid U_3$,
$576\mid U_4$, $14400\mid U_5$ put $C$ on an arithmetic progression), a
mask-arithmetic $2$-adic prefilter to $n=14$ modulo $2^{24}$, and a 64-bit
prefilter for $p=3,5,7,11$ to $n=12$. Survivors are
re-verified with exact rational arithmetic to $n=120$, which also computes the
characteristic roots, the Frobenius obstruction at every point with an integral
exponent difference, the companion $b_n$ ($b_0=0$, $b_1=1$) and its sharp
denominator exponent $k$ (least $k$ with $d_n^kb_n\in\Z$ for $n\le60$,
$d_n=\operatorname{lcm}(1,\dots,n)$), the Ap\'ery limit to $\ge60$ digits, and
the score $\log(1/|\lambda_2|)-k$.

Two validations are on record: the deep scanner returns bit-identical hit lists
to the first scanner on $14$ classes over a common box, and the mixed-class
scanner agrees with an exact brute force on a small box in the class
$(-\tfrac12,0;1,0,0)$, both returning the same $4$ rows --- soundness \emph{and}
completeness \Verified{as stated}.

\subsection{Deduplication}\label{sec:dedup}
A hit is discarded or reclassified by the following tests, in order.
\begin{enumerate}[leftmargin=3.1em]
\item[(D1)] \emph{Repeated characteristic root}, $\Delta=0$: two of the $t_i$
collide.
\item[(D2), (D3)] \emph{Terminating rows} --- three consecutive zeros kill the
sequence (in the class $\mathbf S_1$ this alone removes all of the
$6.5\times10^5$ raw hits at $|a|\le40$) --- and \emph{constant-coefficient rows},
$P,Q,R$ all proportional to $(n+1)^2$ (Proposition~\ref{prop:F5}).
\item[(D4)] \emph{Casoratian degeneracy.} For three solutions of \eqref{eq:four} the
Casoratian $W_n=\det[u^{(i)}_{n-2+k}]$ satisfies $W_{n+1}=R(n)W_n/(n+1)^2$
\Proved, so if $R(n_0)=0$ for an integer $n_0\ge1$ the solution space collapses
to two dimensions from $n_0$ on.
\item[(D5)] \emph{Apparent singularity $\Rightarrow$ disguised three-term row}
(\S\ref{sec:apparent}): the projective local system has only four singular
points, and the row is a gauge or cusp-move image of a three-term row, over
$\Q(\lambda_i)$ in general since the M\"obius map $t\mapsto t/(1-\lambda_it)$ is
defined over that field.
\item[(D6)] \emph{Rescaling}: $u_n\mapsto\mu^nu_n$ acts by
$(a,c,d,f,C)\mapsto(\mu a,\mu c,\mu^2d,\mu^2f,\mu^3C)$; a hit is marked
\textsc{rescaled} when dividing by some $\mu\ge2$ leaves an integral row.
\item[(D7), (D8)] \emph{Pullbacks $t\mapsto t^2$}: these would require $\Rc(-t)=\Rc(t)$,
i.e.\ $a=g=0$, but $g\ne0$, so \emph{there are none in the four-term world}
\Proved. \emph{The Pad\'e $\log$/$\arctan$ families} are three-term and their
four-term images are removed by (D5); exactly two surviving five-point rows have
$k=1$ and are marked Pad\'e-type by the same diagnostic.
\end{enumerate}

Two mechanisms produce (D5) hits and account for essentially all of them: the
\emph{generic cusp move} $s=t/(1-\nu t)$ with $\nu$ \emph{not} a characteristic
root, which sends the ordinary point $t=1/\nu$ to $s=\infty$ and makes that new
fifth point apparent; and a \emph{rational gauge at an ordinary point},
$y\mapsto y/(1-\nu t)$ with $\Rc(1/\nu)\ne0$, which creates a fifth point with
exponents $(0,-1)$ and no logarithm --- and unequal $\rho_i$, so a reducible
cubic, in accordance with Corollary~\ref{cor:galois}.

\begin{proposition}[the generic cusp move, four-term form; \Proved]\label{prop:F4}
Let $(P_0,Q_0)$ be a three-term row in Zagier's class,
$P_0=a_0(n^2+n)+c_0$, $Q_0=d_0n^2$, with integral $u_n$, and let $\nu\in\Z$ be
arbitrary \textup{(}in particular \emph{not} a characteristic root\textup{)}.
Then the signed binomial transform
$v_n=\sum_{m=0}^n\binom nm(-\nu)^{\,n-m}u_m$ satisfies the four-term row of class
$(\rho;M,j_1,j_2)=(0;1,0,1)$ with
\[
P(n)=(a_0-3\nu)(n^2+n)+(c_0-\nu),\quad
Q(n)=(d_0-2a_0\nu+3\nu^2)n^2,\quad
R(n)=-\nu(d_0-a_0\nu+\nu^2)n(n-1),
\]
i.e.\ $f=Q(0)=0$ identically. Its characteristic roots are
$\{\lambda_1-\nu,\lambda_2-\nu,-\nu\}$, and $t=\infty$ is an \emph{apparent}
singularity: the local system is the original four-point one.
\end{proposition}

\begin{proof}
$\sum_nv_nt^n=(1+\nu t)^{-1}y\bigl(t/(1+\nu t)\bigr)$, so the singular points of
the transform are the preimages $t=1/(\lambda_i-\nu)$ of the $t_i$, the preimage
$t=-1/\nu$ of $s=\infty$, and $t=\infty$ --- the image of the \emph{ordinary}
point $s=1/\nu$, carrying only the pole of the gauge factor, hence exponents
$(0,1)$ and no logarithm. Then
$\Rc(t)=\prod_i(1-(\lambda_i-\nu)t)\cdot(1+\nu t)$ gives $a,d,g$, $c=v_1=c_0-\nu$,
and $f$ falls out of $4v_2=P(1)v_1-Q(1)$.
\end{proof}

Proposition~\ref{prop:F4} was checked on Zagier $\mathbf A$--$\mathbf F$ for
$|\nu|\le3$ \Verified{$n\le18$}. For a three-term row of another class the gauge
exponent must be taken in $\{1-r_1,1-r_2\}$; the resulting classes are those with
$\delta_\infty\in\Z$, and every one is populated in the scan. This is the
four-term analogue of the $I_m^*/I_0^*$ family classes of the rigid case
\cite{CLASS}: there the infinite families were the classical Legendre--Pad\'e
rows; here they are cusp-move images, a fifth singular point that is not really
there.

\subsection{The boxes}\label{sec:boxes}
\emph{Box G} (general census): $1\le a\le200$ ($\le140$ for the classes of the
third pass), $|c|\le40$, $|d|\le1200$, $|f|\le80$, $1\le|C|\le120$, integrality
prime-tested to $n=24$ and re-verified exactly to $n=120$.

\emph{Box W} ($|\lambda_2|<1$; complete). A positive score needs
$|\lambda_2|<1$, hence $|\lambda_3|\le|\lambda_2|<1$, and from $a=\sum\lambda_i$,
$d=\sum_{i<j}\lambda_i\lambda_j$, $g=\lambda_1\lambda_2\lambda_3$ this forces
\Proved
\[
|\lambda_1|\le|a|+2,\qquad |d|\le2|a|+5,\qquad |C|=|g|/M^2\le(|a|+2)/M^2 .
\]
Box W takes $|c|\le60$, $|f|\le120$ and lets $|d|,|C|$ run over exactly those
ranges; swept to $a\le300$ for $18$ classes and $a\le400$ for $\mathbf S_0$, it
is \emph{complete} for the question ``is there a row with $|\lambda_2|<1$?''.

\emph{Box PW} (positive-score window; complete). Every census row has $k\ge1$, so
a positive score needs $|\lambda_2|<e^{-1}$, whence \Proved
\[
|\lambda_1|\le|a|+2e^{-1},\qquad |d|\le0.74(|a|+1)+2,\qquad
|C|\le\bigl(0.136(|a|+1)+1\bigr)/M^2 .
\]
Box PW is defined by $1\le a\le1500$, $|c|\le60$, $|f|\le120$ over all $27$
classes and was swept to completion for $a\le800$; inside it the scan is complete
for the question ``is there a row with positive score?'' among rows with
$k\ge1$.

\emph{The deep boxes.} A later run added box DG (the $15$ carrier classes,
$|c|\le200$, $|d|\le2000$, $|f|\le400$, $|C|\le150$), box DW (all $27$
equal-exponent classes, $|c|\le250$, $|f|\le500$, complete in $(d,C)$ for every
$a\le180$, about $17\times$ the $(c,f)$ volume of box W), box MIX (the $20$ mixed
classes, $a\le400$, $|c|\le100$, $|d|\le3000$, $|f|\le200$, $|C|\le400$, sharded
by the rational characteristic root $r$) and box F5 (seven five-term classes,
\S\ref{sec:sixpoint}). The cost is proportional to the product of the five
ranges, so the boxes were sized to budget with the accessory directions $c=u_1$
and $f=Q(0)$ --- where the period lives --- given priority over $a$ and $C$.

\subsection{What the scans produce}\label{sec:scanshape}
Raw hit counts are dominated by degenerate and disguised rows. Two classes make
the point on their own: in $\mathbf S_1$ the box $a\le40$, $|c|\le30$,
$|d|\le600$, $|f|\le50$, $|C|\le80$ contains $646\,297$ parameter tuples passing
integrality mod $p\le24$ to $n=24$ with $\Delta\ne0$, and every one of them has
$u_n=0$ for $n\ge1$; in $\mathbf S_{-1}$ the same box contains $484\,697$, and
every one of them is the constant-coefficient row of Proposition~\ref{prop:F5}.
After (D2)--(D3) both classes are empty.

Over the whole census --- $27$ classes, box G for $21$ of them and a probe box
for the six purely disguised ones --- the filters act as follows
\Verified{boxes as stated}:

\begin{center}\small
\begin{tabular}{lr}
\toprule
stage & rows\\
\midrule
integral mod $p\le24$ to $n=24$, $\Delta\ne0$, non-terminating, not constant-coefficient & $2178$\\
of which exact integrality to $n=120$ fails & $4$\\
of which (D5) an apparent singularity: four-point local system, a disguised three-term row & $2138$ $(98.2\%)$\\
of which (D6) a rescaling $u_n\mapsto\mu^nu_n$ of another integral row & $6$\\
\textbf{genuine five-point rows, primitive} & $\mathbf{30}$\\
\bottomrule
\end{tabular}
\end{center}

Two observations are stable across all boxes. First, the classes producing
$10^5$--$10^6$ raw hits are exactly the Casoratian-degenerate ones; in each the
hits are terminating or constant-coefficient, so after (D2)--(D3) the class is
empty, the one survivor being $(0;1,0,1)$, $R=Cn(n-1)$ --- the class of
Proposition~\ref{prop:F4} --- all of whose rows are disguised. Second, every
class with $\rho\in\Z\setminus\{0\}$ or $\delta_\infty\in\Z_{>0}$ produces a large
family with vanishing Frobenius obstruction throughout, the extreme case being
$(1;2,3,5)$ with \emph{two} apparent points, whose rows are hypergeometric. The
classes with $\rho\notin\Z$, or with $\rho=0$ and $\delta_\infty\notin\Z$, admit
no apparent singularity at all, and they carry the census.

The later deep sweep confirms the picture and adds nothing to it: over boxes DG
and DW (DG $a\le12$, DW complete to $a\le180$) there are $251$ distinct integral
rows, $232$ ($92.4\%$) disguised, $1$ rescaled and $18$ candidates --- \emph{all
$18$ already in the list of $30$}, so zero new equal-exponent five-point rows in
the region covered \Verified{as stated}. The disguise fraction is $92\%$ rather
than $98\%$ only because the deep box is weighted towards the carrier classes.

% ===========================================================================
\section{The sporadic rows and their periods}\label{sec:sporadic}
% ===========================================================================

\subsection{The list of thirty}\label{sec:thirty}
Table~\ref{tab:thirty} lists all $30$ primitive rows with five genuine singular
points produced by boxes G and W, sorted by score. $(\rho;M,j_1,j_2)$ is the
class, $(a,c,d,f,C)$ the parameters of \eqref{eq:class}, $\lambda_i$ the
characteristic roots, $k$ the sharp denominator exponent of the companion,
$\text{score}=\log(1/|\lambda_2|)-k$, and $\Jm$ the outcome of the
$\deg\le18$ $\Jm$-map test.

\begin{table}[htbp]\centering\scriptsize
\setlength{\tabcolsep}{3.5pt}
\caption{The $30$ genuine five-singular-point rows. \Verified{integrality to
$n=120$; $k$ to $n=60$; obstructions at $100$ digits}}\label{tab:thirty}
\begin{tabular}{llllrrrl}
\toprule
class & $(a,c,d,f,C)$ & $\lambda_1,\lambda_2,\lambda_3$ & $k$ & score & $\Ic_1$ & $\Ic_2$ & $\Jm$\\
\midrule
$(0;2,1,1)$ & $(154,42,-128,-12,8)$ & $154.8,\,-0.414\pm0.1878i$ & $2$ & $-1.2117$ & $456533/4$ & $-5929/32$ & no\\
$(0;2,1,1)$ & $(87,24,-152,-16,16)$ & $88.72,\,-1,\,-0.7214$ & $2$ & $-2.0000$ & $658503/64$ & $-7569/152$ & no\\
$(\tfrac12;4,5,5)$ & $(72,6,-192,-60,8)$ & $74.6,\,-1.298\mp0.1732i$ & $2$ & $-2.2700$ & $2916$ & $-27$ & no\\
$(0;3,-1,4)$ & $(22,6,-72,12,36)$ & $25.35,\,-1.673\pm3.16i$ & $1$ & $-2.2741$ & $2662/81$ & $-121/18$ & no\\
$(0;2,1,1)$ & $(22,6,-96,-12,24)$ & $25.86,\,-2,\,-1.856$ & $2$ & $-2.6931$ & $1331/12$ & $-121/24$ & no\\
$(0;2,1,1)$ & $(11,3,-73,-8,28)$ & $16,\,-2.5\pm0.866i$ & $2$ & $-2.9730$ & $1331/112$ & $-121/73$ & no\\
$(-\tfrac12;4,-1,-1)$ & $(16,10,64,20,4)$ & $10.47,\,4,\,1.528$ & $2$ & $-3.3863$ & $64$ & $4$ & no\\
$(\tfrac12;4,5,5)$ & $(16,2,64,20,4)$ & $10.47,\,4,\,1.528$ & $2$ & $-3.3863$ & $64$ & $4$ & no\\
$(0;2,1,1)$ & $(60,18,240,12,16)$ & $55.71,\,4,\,0.2872$ & $2$ & $-3.3863$ & $3375$ & $15$ & no\\
$(0;2,1,1)$ & $(20,6,-80,-4,-16)$ & $23.31,\,-4,\,0.6863$ & $2$ & $-3.3863$ & $-125$ & $-5$ & no\\
$(0;2,1,1)$ & $(3,0,-104,-16,80)$ & $12.94,\,-5,\,-4.944$ & $2$ & $-3.6094$ & $27/320$ & $-9/104$ & no\\
$(0;3,1,2)$ & $(11,4,37,3,3)$ & $5\pm1.414i,\,1$ & $2$ & $-3.6479$ & $1331/27$ & $121/37$ & \textbf{yes}, $(4,16)$\\
$(0;2,1,1)$ & $(16,6,96,12,64)$ & $8,\,4\pm4i$ & $2$ & $-3.7329$ & $16$ & $8/3$ & no\\
$(0;2,1,1)$ & $(16,6,76,8,24)$ & $8,\,6,\,2$ & $2$ & $-3.7918$ & $128/3$ & $64/19$ & no\\
$(0;2,1,1)$ & $(16,6,96,12,48)$ & $6\pm3.464i,\,4$ & $2$ & $-3.9356$ & $64/3$ & $8/3$ & no\\
$(0;4,1,3)$ & $(14,5,97,8,12)$ & $5.5\pm5.809i,\,3$ & $2$ & $-4.0794$ & $343/24$ & $196/97$ & no\\
$(0;2,1,1)$ & $(16,6,16,-4,-96)$ & $12,\,8,\,-4$ & $2$ & $-4.0794$ & $-32/3$ & $16$ & no\\
$(0;6,-1,7)$ & $(13,4,432,-24,-48)$ & $8.262\pm20.54i,\,-3.525$ & $1$ & $-4.0975$ & $-2197/1728$ & $169/432$ & no\\
$(0;2,1,1)$ & $(6,2,64,4,-40)$ & $4\pm8i,\,-2$ & $2$ & $-4.1910$ & $-27/20$ & $9/16$ & no\\
$(0;3,1,2)$ & $(3,0,-216,-24,108)$ & $18,\,-9,\,-6$ & $2$ & $-4.1972$ & $1/36$ & $-1/24$ & no\\
$(0;3,1,2)$ & $(33,12,324,36,108)$ & $18,\,9,\,6$ & $2$ & $-4.1972$ & $1331/36$ & $121/36$ & no\\
$(0;3,-1,4)$ & $(3,0,-216,48,108)$ & $18,\,-9,\,-6$ & $2$ & $-4.1972$ & $1/36$ & $-1/24$ & no\\
$(0;2,1,1)$ & $(17,6,112,8,24)$ & $8\pm5.657i,\,1$ & $2$ & $-4.2822$ & $4913/96$ & $289/112$ & no\\
$(0;4,1,3)$ & $(28,10,208,12,16)$ & $16,\,10.47,\,1.528$ & $2$ & $-4.3487$ & $343/4$ & $49/13$ & no\\
$(0;6,1,5)$ & $(17,6,56,0,-12)$ & $10.21\pm4.658i,\,-3.428$ & $2$ & $-4.4182$ & $-4913/432$ & $289/56$ & no\\
$(0;2,1,1)$ & $(32,12,272,16,64)$ & $16,\,14.93,\,1.072$ & $2$ & $-4.7033$ & $128$ & $64/17$ & no\\
$(0;2,1,1)$ & $(42,15,441,24,100)$ & $25,\,16,\,1$ & $2$ & $-4.7726$ & $9261/50$ & $4$ & no\\
$(-\tfrac13;1,0,0)$ & $(42,22,480,64,-64)$ & $21.07\pm6.464i,\,-0.1318$ & $2$ & $-5.0926$ & $-9261/8$ & $147/40$ & no\\
$(\tfrac13;1,1,1)$ & $(42,8,480,64,-64)$ & $21.07\pm6.464i,\,-0.1318$ & $2$ & $-5.0926$ & $-9261/8$ & $147/40$ & no\\
$(0;2,1,1)$ & $(64,24,1030,26,48)$ & $32,\,31.81,\,0.1886$ & $2$ & $-5.4598$ & $4096/3$ & $2048/515$ & no\\
\bottomrule
\end{tabular}
\end{table}

Four notes.

\begin{enumerate}[label=\textup{\arabic*.},leftmargin=2.1em]
\item $k=2$ for $28$ of the $30$ --- the four-term echo of the free integration
phenomenon of the rigid case. The two exceptions have $k=1$, lie in the classes
with $\delta_\infty>1$, and are Pad\'e constructions rather than modular objects
by the $k=1$ diagnostic.
\item Exactly one row is an elliptic surface: the $K3$ of \S\ref{sec:k3}. The
other $29$ fail the $\Jm$-map test over the whole $\deg\le18$ window --- the
four-term half of the non-congruence picture. A \emph{positive} $\Jm$ answer is a
certificate; a \emph{negative} one means only ``no fit with $h\le18$, $\deg\le18$
and $\gamma$ in the $760$-element list'' \Verified{as stated}.
\item \emph{Quadratic-twist pairs.} $(16,10,64,20,4)$ and $(16,2,64,20,4)$, and
likewise $(42,22,480,64,-64)$ and $(42,8,480,64,-64)$, have the same $(a,d,f,C)$
and differ only in the accessory $c$, hence have the same five singular points;
in each pair the classes are exchanged by the twist
$\prod_i(1-\lambda_it)^{\mp1/2}$, so they are the two half-integral gauges of one
projective local system. \emph{Their Ap\'ery limits differ.}
\item The best score is $-1.21$, attained by $(154,42,-128,-12,8)$ in
$\mathbf S_0$, whose two small roots are the complex pair $-0.4140\pm0.1878i$;
its Ap\'ery limit $\xi=0.02358580301301982283\dots$ is unidentified. Adding a
fifth \emph{genuine} singular point to a known four-point configuration is a
recurring motif: $(87,24,-152,-16,16)$ has
$\chi=(\lambda+1)(\lambda^2-88\lambda-64)$, the quadratic being the level-$5$
Fricke row's, and $(20,6,-80,-4,-16)$ has $(\lambda+4)(\lambda^2-24\lambda+16)$,
the quadratic being Beukers' row's.
\end{enumerate}

\subsection{The periods of the thirty: a negative result}\label{sec:negper}
For each row the companion $b_n$ satisfies $LB=t$, and $\xi=\lim b_n/a_n$ exists
exactly when the largest characteristic root is real and simple, with error
$O((\lambda_2/\lambda_1)^n)$. Nine of the thirty have a complex conjugate
dominant pair and therefore \emph{no archimedean Ap\'ery limit} --- in
particular the $K3$ row, whose dominant roots are $5\pm i\sqrt2$.

The $21$ limits that do exist were computed to $\ge60$ digits and put through two
batteries: an integer-relation search over $103$ classical weight-two constants
($\zeta(2)$; $L(2,\chi_D)$ for $|D|\le24$; $\log^2m$, $\pi\log m$ and their
quadratic-unit variants; $\pi\arctan\frac1m$, $\arctan^2\frac1m$;
$\Gamma(1/4)^4/\pi$, $\Gamma(1/3)^6/\pi^2$; $\mathrm{Li}_2$ at twenty small
rational arguments), and, for each row separately, the weight-two polylogarithm
space attached to \emph{its own} singular points, built from the ratios
$\lambda_i/\lambda_j$ and $1-\lambda_i/\lambda_j$ and reduced to a greedy
independent subset first (without that reduction the classical dilogarithm
identities make the basis dependent and the search returns them instead of
anything about $\xi$).

\emph{All $21$ limits came back unidentified}, in both batteries, at $60$ digits
with coefficient heights up to $10^6$ (single constant) and $10^5$ (pairs), and
in polylogarithm spaces of dimension $7$--$11$ \Verified{stated bases and
heights} --- a negative result, but sharp, and consistent with the geometry: a
five-point configuration is not a modular curve.

\subsection{The complex fold constants}\label{sec:folds}
For the nine rows with no archimedean limit the natural invariants are the
\emph{fold constants}: continue $(A,A')$ and $(B,B')$ from inside
$|t|<\min|t_i|$ to a singular point $t_i$, match to the local Frobenius basis
(logarithmic for $\rho\in\Z$, $\{1+\cdots,\ s^\rho(1+\cdots)\}$ for
$\rho\notin\Z$), and set
\[
\xi(t_i)=\frac{\text{coefficient of the singular basis element in }B}
              {\text{coefficient of the singular basis element in }A}.
\]
The machinery self-tests on the $K3$ row, reproducing its three published
constants to $210$--$211$ digits, and its non-integer-$\rho$ branch was validated
on the two $\rho=\pm\tfrac13$ rows. All nine rows succeed: $27$ fold constants,
each self-agreeing to $154$--$156$ digits between working precisions $130$ and
$210$ \Verified{as stated}. In each of the nine the three finite singular points
are one real point and a conjugate pair, so each row contributes three real
periods $\xi(t_1)$, $\xi_++\xi_-$, $(\xi_+-\xi_-)/i$; and $\xi(t_1)=\xi_++\xi_-$
only for the $K3$ row, so the other eight contribute two independent real periods
each.

Put through the $145$-constant battery of \S\ref{sec:battery}, the only hits are
the $K3$ row's, returning the expected $L(g,2)$ multiples at $124$--$125$ digits
--- a strong positive control. \emph{Every one of the $24$ genuinely new real
periods is unidentified} \Verified{as stated}.

One cross-row coincidence stands out: $\re\xi(t_2)=0.14639462825244069157\dots$
of the row $(0;4,1,3)$, $(14,5,97,8,12)$ equals the archimedean Ap\'ery limit of
the \emph{different} row $(0;2,1,1)$, $(42,15,441,24,100)$ to all $70$ stored
digits. The two operators have different singular points ---
$\lambda=\{3,\tfrac{11\pm i\sqrt{115}}2\}$ against $\{1,16,25\}$ --- so this is
neither a twist nor a rescaling \Numerical{$70$ digits}. Unexplained; see
Conjecture~\ref{conj:D6}.

\subsection{The elliptic $K3$ row}\label{sec:k3}
The row
\begin{equation}\label{eq:k3row}
(n+1)^2u_{n+1}=(11n^2+11n+4)u_n-(37n^2+3)u_{n-1}+3(3n-1)(3n-2)u_{n-2},
\end{equation}
$u_n=1,4,16,64,250,928,3136,8704,11866,-79400,\dots$, lies in the class
$(0;3,1,2)$, i.e.\ $(\rho;\delta_\infty)=(0;\tfrac13)$: exponents $(0,0)$ at
$t=0$ and at each $t_i$, and $(\tfrac53,\tfrac43)$ at $\infty$. Its
characteristic polynomial is
$\lambda^3-11\lambda^2+37\lambda-27=(\lambda-1)(\lambda^2-10\lambda+27)$, so the
singular points are $0,\ 1,\ \tfrac{5\pm i\sqrt2}{27},\ \infty$; integrality is
verified exactly to $n=300$ ($u_{300}$ has $213$ decimal digits) and $k=2$
\Verified{as stated}. The two dominant characteristic roots are the conjugate
pair $5\pm i\sqrt2$, so there is no archimedean Ap\'ery limit.

The $\Jm$-map test is positive, returning $h=4$, $\deg\Jm=16$, $\gamma=-1$,
certified: $\Jm=U/V$ with $\deg U=15$, $\deg V=16$ fitted from $38$ coefficients
of the $q$-expansion and verified on $73$ further ones \Verified{as stated}.
Factoring $V$, $U$ and $U-1728V$ gives the fibres $I_4$ at $t=0$, $I_6$ at $t=1$,
$I_3,I_3$ at $t=\tfrac{5\pm i\sqrt2}{27}$ and $IV^*$ at $\infty$, with
$\Jm^{-1}(\infty):4+6+3+3=16$, $\Jm^{-1}(0):3\cdot5+1=16$,
$\Jm^{-1}(1728):2\cdot8=16$ and $\sum_P(e_P-1)=30=2\cdot16-2$: $\Jm$ is a
\emph{Belyi map}, so there are no further singular fibres, and the projective
monodromy is the genus-zero group of signature $(0;3;\ \text{cusps }3,3,4,6)$ and
index $16$ --- the $(c,e_2,e_3)=(4,0,1)$ line of \S\ref{sec:RH}, at the
Riemann--Hurwitz bound $6\cdot4+4\cdot1-12=16$ with equality.

Since $\sum_{I_n}n_i=16>12$ the surface is not rational; the Euler numbers give
$\sum e=4+6+3+3+8=24$, so $\chi=2$: an \emph{elliptic $K3$ surface over
$\PP^1_t$ with exactly five singular fibres $I_4I_6I_3I_3IV^*$}, defined over
$\Q$ with the two $I_3$'s conjugate over $\Q(\sqrt{-2})$.

\begin{theorem}[\cite{K3PAPER,K3P}; stated here, not reproved]\label{thm:k3}
For the row \eqref{eq:k3row}: the local system is $R^1\pi_*\Q$ of an elliptic
$K3$ surface with the fibres above, given by an explicit Weierstrass model
\Proved; $\rho(\mathcal X)=20$, $\mathrm{MW}$ has rank $0$ and torsion exactly
$\Z/3$ \Proved; the transcendental lattice is $\operatorname{diag}(6,12)$,
$\det T=72$, with CM by $\Q(\sqrt{-2})$ \Proved{} modulo the discriminant-form
bookkeeping; the transcendental motive is that of the weight-$3$ CM newform
$g=\texttt{32.3.d.a}\in S_3(\Gamma_0(32),\chi_{-8})$ \Verified{$a_p$ for all $76$
primes $5\le p\le397$, no free parameter, Livn\'e rigidity then identifying the
motive \cite{Livne}}; and the fold constants are
\[
\xi\Bigl(\tfrac{5\mp i\sqrt2}{27}\Bigr)=\frac{\sqrt2\mp i}{3}L(g,2),
\qquad
\xi(1)=\frac{2\sqrt2}{3}L(g,2)=\xi_++\xi_-,
\qquad
L(g,2)=\frac{\pi}{32}\frac{\Gamma(\tfrac18)\Gamma(\tfrac38)}{\Gamma(\tfrac58)\Gamma(\tfrac78)}
\]
\Verified{$\ge165$ digits, two independent computations}. $2$-adically
$b_n/a_n\to0$ at logarithmic rate and there is no $3$-adic limit
\Verified{as stated in \cite{K3P}}.
\end{theorem}

This is the structural pay-off of the fifth point: \emph{the four-term row is the
recurrence-level home of the elliptic $K3$s with five singular fibres, just as
the three-term row is the home of the rational elliptic surfaces with four.} The
integrality of \eqref{eq:k3row} is nonetheless purely experimental \Open --- for
three-term rows integrality is explained by modular parametrisation, and here we
do not know whether the index-$16$ monodromy group is congruence.

\subsection{The mixed-exponent rows: Catalan and $L$-values}\label{sec:mixed}
The mixed classes of \S\ref{sec:mixdict} were scanned next, and they carry
everything the equal-exponent scan does not. Six genuine five-point rows with
identified classical periods appear, all in the single class
\[
(\rho_p,\rho_r;M,j_1,j_2)=(-\tfrac12,0;1,0,0),\qquad R(n)=Cn^2,
\]
i.e.\ exponent differences $(0,\tfrac12,\tfrac12)$ at the three finite points
and $\delta_\infty=0$: one $I_n$ fibre and two points whose projective monodromy
has order $2$ (Kodaira $III$/$III^*$), with $I_n$'s at $0$ and $\infty$.

\begin{table}[htbp]\centering\footnotesize
\setlength{\tabcolsep}{4pt}
\caption{The six mixed-exponent rows with classical Ap\'ery limits.
\Verified{integrality to $n=200$; exponents to $40$ digits; $k$ to $n=80$;
$\xi$ to $139$--$141$ digits}}\label{tab:six}
\begin{tabular}{clllllrl}
\toprule
\# & $P(n)$ & $Q(n)$ & $R(n)$ & $u_n$ & $\lambda$ & score & $\xi=\lim b_n/a_n$\\
\midrule
1 & $16n^2{+}20n{+}8$ & $48n^2{+}16n$ & $-128n^2$ & $1,8,72,640,\dots$ & $4{\pm}4\sqrt2,\,8$ & $-4.0794$ & $\tfrac14G$\\
2 & $14n^2{+}20n{+}8$ & $28n^2{+}16n{+}4$ & $8n^2$ & $1,8,72,704,\dots$ & $6{\pm}4\sqrt2,\,2$ & $-2.6931$ & $\tfrac12G-\tfrac3{16}\zeta(2)$\\
3 & $6n^2{+}10n{+}4$ & $-32n^2{-}24n{-}8$ & $32n^2$ & $1,4,36,288,\dots$ & $4{\pm}4\sqrt2,\,-2$ & $-2.6931$ & $\tfrac38\zeta(2)-\tfrac12G$\\
4 & $16n^2{+}20n{+}8$ & $68n^2{+}36n{+}8$ & $32n^2$ & $1,8,60,448,\dots$ & $8,\,4{\pm}2\sqrt3$ & $-4.0101$ & $\tfrac1{32}(2\zeta(2)+15L_{-3})$\\
5 & $17n^2{+}25n{+}10$ & $32n^2{+}24n{+}8$ & $16n^2$ & $1,10,114,1424,\dots$ & $8{\pm}4\sqrt3,\,1$ & $-2.0693$ & $\tfrac1{16}(15L_{-3}-6\zeta(2))$\\
6 & $13n^2{+}20n{+}8$ & $-13n^2{-}6n{-}1$ & $-n^2$ & $1,8,87,1024,\dots$ & $7{\pm}4\sqrt3,\,-1$ & $-2.0000$ & $\tfrac18(2\zeta(2)-3L_{-3})$\\
\bottomrule
\end{tabular}
\end{table}

Here $G=L(2,\chi_{-4})$, $L_{-3}=L(2,\chi_{-3})$, all six have $k=2$, and the
class parameters $(a,c,d,f,C)_r$ are $(16,8,48,0,-128)_{8}$,
$(14,8,28,4,8)_{2}$, $(6,4,-32,-8,32)_{-2}$, $(16,8,68,8,32)_{8}$,
$(17,10,32,8,16)_{1}$, $(13,8,-13,-1,-1)_{-1}$.
Independently of the scan pipeline the coefficients were rebuilt from
Theorem~\ref{thm:D3} and, for each row: $u_n\in\Z$ by exact rational arithmetic
to $n=200$; the exponents recomputed from the dictionary come out as
$(0,-\tfrac12,-\tfrac12)$ to $40$ digits with cubic discriminants
$2048,32768,32768,768,192,49152$, all non-zero, so the five singular points are
distinct; \emph{no apparent singularity is possible}, since $\rho_r=0$ and
$\delta_\infty=0$ force logarithms and the two $III$ points have monodromy of
order $2$ in $\PSL_2$; $k=2$ to $n=80$; $\xi$ recomputed by forward recursion at
working precision $380$ agrees to $139$--$141$ digits; and \emph{the sequences
satisfy no three-term recurrence} --- the nine-parameter ansatz
$L(n)u_{n+1}=P(n)u_n-Q(n)u_{n-1}$ with $L,P,Q$ quadratic has nullspace dimension
$0$ over $21$ equations, while the twelve-parameter four-term ansatz has
nullspace dimension exactly $1$. \Verified{as stated}

The finite singular points of rows 1--3 are irrational and conjugate over
$\Q(\sqrt2)$, the field of the level-$8$ and level-$16$ Catalan hosts; rows 4--6
are the same phenomenon over $\Q(\sqrt3)$. Rows 2 and 3 realise the mixed class
$\tfrac12G\mp\tfrac3{16}\zeta(2)$, $\tfrac38\zeta(2)-\tfrac12G$ on a five-point
host. The best Catalan score is $-2.6931=\log\tfrac12-2$, i.e.\ $+0.693$ nats
better than Zagier $\mathbf E$'s $-3.3863$, and row 6 reaches $-2.0000$ with
$|\lambda_2|=1$ exactly. For rows $1,2,3,5,6$ the singular point nearest the
origin --- the one that governs $\lim b_n/a_n$ --- is one of the two $III$
points, an order-$2$ orbifold point rather than a cusp; for row 4 it is the cusp
$t=\tfrac18$. Every one of the six has $|\lambda_2|\ge1$, so none is a decayer,
as Theorem~\ref{thm:D1} forces.

\subsection{The whole period matrix is classical}\label{sec:matrix}
The Ap\'ery limit is only the fold constant at the nearest singular point. The
constants at all three were computed and identified, real and imaginary parts,
at $125$--$129$ digits \Verified{as stated}. Ordering $t_1,t_2,t_3$ by
decreasing $|t|$, so that $\xi(t_3)$ is the archimedean Ap\'ery limit:

\begin{center}\footnotesize
\begin{tabular}{llll}
\toprule
row & $\xi(t_1)$ & $\xi(t_2)$ & $\xi(t_3)$\\
\midrule
$(16,8,48,0,-128)_{8}$ & $\tfrac18(2G-3\zeta(2))$ & $\tfrac1{32}(8G+3\zeta(2))+\tfrac3{32}\zeta(2)i$ & $\tfrac14G$\\
$(14,8,28,4,8)_{2}$ & $\tfrac1{80}(40G-39\zeta(2))+\tfrac9{10}\zeta(2)i$ & $\tfrac1{16}(8G+3\zeta(2))+\tfrac38\zeta(2)i$ & $\tfrac1{16}(8G-3\zeta(2))$\\
$(16,8,68,8,32)_{8}$ & $\tfrac1{32}(15L_{-3}-12\zeta(2))+\tfrac5{96}\pi^2\sqrt3\,i$ & $\tfrac{15}{32}L_{-3}+\tfrac1{96}\pi^2\sqrt3\,i$ & $\tfrac1{32}(2\zeta(2)+15L_{-3})$\\
$(8,4,32,8,64)_{4}$ & $\tfrac1{32}(6\zeta(2)+15L_{-3})$ & $\tfrac1{32}(15L_{-3}-6\zeta(2))+\tfrac1{48}\pi^2\sqrt3\,i$ & $\overline{\xi(t_2)}$\\
\bottomrule
\end{tabular}
\end{center}

Two consequences. First, the period lattice is spanned by two classical
constants: for the $\Q(\sqrt2)$ family every fold period lies in
$\Q G+\Q\zeta(2)$, for the $\Q(\sqrt3)$ family in $\Q L_{-3}+\Q(\sqrt3)\zeta(2)$
(note $\pi^2\sqrt3=6\sqrt3\zeta(2)$). These are \emph{not} generic five-point
periods --- all $21$ equal-exponent Ap\'ery limits and all $24$ new
equal-exponent fold periods are unidentified over the same battery --- so the
mixed class is arithmetically special. Second, the row $(8,4,32,8,64)_{4}$ has no
archimedean Ap\'ery limit at all, its dominant roots being the complex pair
$2\pm2\sqrt3\,i$, and yet all three of its fold periods are classical: the
phenomenon belongs to the local system, not to the recurrence's convergence.

\subsection{The gauge partners and $\Gamma(\tfrac14)^4$}\label{sec:gauge}
The class $(-\tfrac12,0;1,0,0)$ has the gauge partner $(+\tfrac12,0;1,1,1)$ ---
the same projective local system with $\rho\mapsto-\rho$, the $\pm\tfrac12$ twist
pairing of \S\ref{sec:thirty} note~3. It has $R(n)=C(n-1)^2$, so $R(1)=0$: a
Casoratian degeneracy in the sense of (D4), which is why it was initially left
out of the class list. Its integral rows have the \emph{same} $(a,d,f,C)$ as the
Catalan rows with a different accessory $c$:
\[
\begin{array}{llll}
(16,4,48,0,-128)_{8}:& 1,4,24,160,1136,8384,\dots\\
(14,2,28,4,8)_{2}:& 1,2,8,48,356,2952,\dots\\
(6,0,-32,-8,32)_{-2}:& 1,0,4,16,116,800,\dots
\end{array}
\]
and \emph{all three} have the same Ap\'ery limit,
\[
\xi=\frac{\Gamma(\tfrac14)^4}{64\pi}
 =0.85939822725254660343626197247631964973760705647739511203431098644200041\dots,
\]
\Verified{integrality to $n=200$; $\xi$ matched to $195$, $241$, $241$ digits}.
Three different rows, three different singular configurations, one and the same
period --- in sharp contrast to the $-\tfrac12$ gauge, where the three partners
give three different constants.

The pairing is the point: $\Gamma(\tfrac14)^4/\pi$ is the period of the CM
elliptic curve $y^2=x^3-x$ --- the algebraic-period side of the Gaussian world
--- and $G=L(2,\chi_{-4})$ is its $L$-value side, and the two half-integral
gauges of one five-point projective local system over $\Q(\sqrt2)$ produce one
each; $\sqrt2$ is the ramification of $\Q(i)/\Q$ made visible in the singular
points.

\subsection{The complete census of the two productive classes, and $\xi^*$}
\label{sec:xistar}
The pair of gauge-partner classes was swept to completion over
\[
|r|\le60,\quad |a|\le200,\quad |c|\le200,\quad |d|\le2000,\quad |f|\le400,
\quad |C|\le600
\]
($69$ distinct integral rows, of which $30$ rescaled, $14$ not integral at
$n=120$, and $20$ candidates). \emph{Each class contains exactly nine primitive
genuine five-point rows}, and they correspond one-to-one under the twist
$\rho\mapsto-\rho$ --- same $(a,d,f,C)$ and same $r$, different accessory $c$
\Verified{box as stated}:

\begin{center}\footnotesize
\begin{tabular}{rlllll}
\toprule
$r$ & $(a,d,f,C)$ & $\lambda$ & score & $\xi$ in gauge $\rho_p=-\tfrac12$ & $\xi$ in gauge $\rho_p=+\tfrac12$\\
\midrule
$8$ & $(16,48,0,-128)$ & $4{\pm}4\sqrt2,\,8$ & $-4.0794$ & $\tfrac14G$ & $\Gamma(\tfrac14)^4/(64\pi)$\\
$2$ & $(14,28,4,8)$ & $6{\pm}4\sqrt2,\,2$ & $-2.6931$ & $\tfrac12G-\tfrac3{16}\zeta(2)$ & $\Gamma(\tfrac14)^4/(64\pi)$\\
$-2$ & $(6,-32,-8,32)$ & $4{\pm}4\sqrt2,\,-2$ & $-2.6931$ & $\tfrac38\zeta(2)-\tfrac12G$ & $\Gamma(\tfrac14)^4/(64\pi)$\\
$8$ & $(16,68,8,32)$ & $8,\,4{\pm}2\sqrt3$ & $-4.0101$ & $\tfrac1{32}(2\zeta(2){+}15L_{-3})$ & $\xi^*$\\
$1$ & $(17,32,8,16)$ & $8{\pm}4\sqrt3,\,1$ & $-2.0693$ & $\tfrac1{16}(15L_{-3}{-}6\zeta(2))$ & $\xi^*$\\
$-1$ & $(13,-13,-1,-1)$ & $7{\pm}4\sqrt3,\,-1$ & $-2.0000$ & $\tfrac18(2\zeta(2){-}3L_{-3})$ & $\xi^*$\\
$4$ & $(8,32,8,64)$ & $4,\,2{\pm}2\sqrt3\,i$ & $-3.3863$ & complex pair, no limit & complex pair, no limit\\
$3$ & $(17,123,33,243)$ & $7{\pm}4\sqrt2\,i,\,3$ & $-4.1972$ & complex pair, no limit & complex pair, no limit\\
$-3$ & $(5,24,12,-144)$ & $4{\pm}4\sqrt2\,i,\,-3$ & $-3.9356$ & complex pair, no limit & complex pair, no limit\\
\bottomrule
\end{tabular}
\end{center}

Every row with an archimedean limit in the $-\tfrac12$ gauge is identified; in
the $+\tfrac12$ gauge three give $\Gamma(\tfrac14)^4/(64\pi)$ and three the same
\emph{unidentified} constant
\[
\xi^*=0.7372929961855962401764261978022936979609191432765291852228688851129961062186382\dots
\]
shared by those three rows, by a fourth four-term row in a \emph{different}
mixed class ($(0,\tfrac12;2,1,2)$, $r=12$, $(a,c,d,f,C)=(24,6,192,24,144)$),
\emph{and} by the five-term six-point row of \S\ref{sec:sixpoint} whose
characteristic quartic $(\lambda-4)(\lambda-2)(\lambda+2)(\lambda-1)$ has all
roots rational --- two computations with recurrences of different orders and
different singular sets, agreeing to $319$ digits \Verified{$319$ agreeing
digits}. $\xi^*$ is unidentified by the $145$-constant battery (heights $10^8$
for single constants, $10^5$ for pairs, \texttt{algdep} to degree $6$), and a
hand search against $\Gamma(\tfrac13)^6/\pi^2$, $\Gamma(\tfrac14)^4/\pi$ and
rational multiples returns only height-$10^{55}$--$10^{85}$ noise. Five rows, in
three classes and two worlds, over two fields, with one common transcendental:
identifying $\xi^*$ is the sharpest single question this work leaves behind
\Open.

\subsection{The battery}\label{sec:battery}
The identifications above and the negative results of \S\ref{sec:negper} and
\S\ref{sec:folds} use a Catalan-focused battery of $145$ constants after
value-deduplication: $G$ and six algebraic multiples of it; $\zeta(2)$, $\pi^2$,
$\pi^2\sqrt D^{\pm1}$ for $D\mid24$; $L(2,\chi_D)$ for all $17$ fundamental
discriminants $|D|\le24$; $L(3,\chi_D)$ for $|D|\le8$, $\zeta(3)$ and
$\zeta(3)/\pi^{1,2}$; $\mathrm{Ti}_2$ at $11$ arguments; the $K3$ row's $L(g,2)$
and four twists; seven $\Gamma$/CM-period products of levels $3,4,6,8,12$; and
$\log^2m$, $\pi\log m$, $28$ products $\log m_1\log m_2$, five
real-quadratic-unit pairs, $\pi\arctan\frac1m$ and $\arctan^2\frac1m$. Constants
are built at $250$ digits, the search is run at $60$, and \emph{every} candidate
relation is verified by reconstructing $\xi$ from it with a $55$-digit threshold;
the pair test uses a $24$-element sublist greedily reduced to an independent
basis of $22$. The battery recovers $\tfrac73G$, $\tfrac52-\tfrac37\pi G$,
$\tfrac25G-\tfrac34\zeta(2)+\tfrac16$ (also from a $62$-digit truncated input),
$\sqrt5$, $\tfrac{2\sqrt2}3L(g,2)$ and $\tfrac12G$ from Zagier $\mathbf E$: it is
not silently blind \Verified{as stated}.

% ===========================================================================
\section{Structure theorems}\label{sec:structure}
% ===========================================================================

The census results above are computations over boxes. Three statements explain,
rather than merely record, part of what they show.

\begin{theorem}[\Proved]\label{thm:D1}
Let \eqref{eq:four} be an integral four-term row with $g=CM^2\ne0$ and
characteristic roots $|\lambda_1|\ge|\lambda_2|\ge|\lambda_3|$. If
$|\lambda_2|<1$ then $\chi(\lambda)=\lambda^3-a\lambda^2+d\lambda-g$ is
irreducible over $\Q$.
\end{theorem}

\begin{proof}
Suppose $\chi$ has a rational root. It is monic with integer coefficients, so
that root is an integer $r$, and $r\ne0$ because $\chi(0)=-g\ne0$. If
$r\in\{\lambda_2,\lambda_3\}$ then $|r|\le|\lambda_2|<1$ with $r$ a non-zero
integer, which is impossible. Hence $r=\lambda_1$ and, by Gauss's lemma,
$\chi=(\lambda-r)(\lambda^2-\sigma\lambda+\pi)$ with $\sigma,\pi\in\Z$. Then
$\pi=\lambda_2\lambda_3$ and $|\pi|\le|\lambda_2|^2<1$, so $\pi=0$ and $g=r\pi=0$
--- a contradiction.
\end{proof}

\begin{corollary}[\Proved]\label{cor:D11}
No class with $\rho_1\ne\rho_2=\rho_3$ --- in particular none of the twenty
mixed-exponent classes of \S\ref{sec:mixdict} --- contains a row with
$|\lambda_2|<1$.
\end{corollary}

\begin{proof}
Corollary~\ref{cor:galois} makes unequal exponents force a reducible cubic, and
Theorem~\ref{thm:D1} forbids that when $|\lambda_2|<1$.
\end{proof}

\begin{corollary}[\Proved]\label{cor:D12}
Every row with $|\lambda_2|<1$ has equal exponents at all three finite points, so
its three finite fibres are Galois-conjugate and of one Kodaira type.
\end{corollary}

This explains, rather than merely records, why the $|\lambda_2|<1$ window is
populated only by the equal-exponent classes.

\begin{theorem}[\Proved]\label{thm:D2}
If moreover $\lambda_2\in\R$ \textup{(}a \emph{real decayer}\textup{)} then
$\chi$ is irreducible and \emph{totally real},
$\Delta=18adg-4a^3g+a^2d^2-4d^3-27g^2>0$, and $|g|<|\lambda_1|\le|a|+2$.
\end{theorem}

\begin{proof}
An irreducible cubic over $\Q$ has either three real roots or one real root and a
complex-conjugate pair. In the second case the pair has equal modulus; if the
real root were $\lambda_2$ then $|\lambda_1|=|\lambda_3|$, contradicting
$|\lambda_2|\ge|\lambda_3|$ and $|\lambda_2|<|\lambda_1|$. So all three roots are
real and $\Delta>0$; finally $|g|=|\lambda_1||\lambda_2\lambda_3|<|\lambda_1|$.
\end{proof}

Both statements were run against the $30$ rows of Table~\ref{tab:thirty} and
against all $37\,088$ integer cubics with $|\lambda_2|<1$ in $|a|\le24$,
$|d|,|g|\le70$: $0$ reducible \Verified{as stated}. Exactly one census row has
$|\lambda_2|<1$, namely $(154,42,-128,-12,8)$ with
$|\lambda_2|=0.45462\ldots$, and its cubic is irreducible with
$\Delta=-81\,920\,000<0$, so its two small roots are a complex pair and it is
\emph{not} a real decayer. The only other row with $|\lambda_2|\le1$,
$(87,24,-152,-16,16)$, has $|\lambda_2|=1$ exactly and a \emph{reducible} cubic
$(\lambda+1)(\lambda^2-88\lambda-64)$ --- precisely the boundary case
Theorem~\ref{thm:D1}'s proof leaves open. Consequently, in the four-term world a
real decayer must live in an equal-exponent class with an irreducible totally
real cubic.

The empirical side agrees. Over the whole extended window (all $27$
equal-exponent classes, $|c|\le250$, $|f|\le500$, complete in $(d,C)$ for every
$a\le180$) there are exactly three rows with $|\lambda_2|<1$ \Verified{as
stated}: the disguised $I_0^*$ Pad\'e row $(144,32,96,-32,4)$ of class
$(0;2,-1,3)$ ($|\lambda_2|=0.35064$ real, $k=1$, score $+0.0480$), the disguised
$(36,8,-48,16,4)$ of the same class ($0.65496$, non-real, $k=1$, $-0.5768$), and
the genuine $(154,42,-128,-12,8)$ of $(0;2,1,1)$ ($0.45462$, non-real, $k=2$,
$-1.2117$). Over all $20$ mixed-class candidates the smallest $|\lambda_2|$ is
\emph{exactly} $1$, never below --- Corollary~\ref{cor:D11} in action.

The third structural statement concerns what the disguise costs. The only
integral three-term row in Zagier's sense whose Ap\'ery limit involves Catalan's
constant is Zagier $\mathbf E$,
$(n+1)^2u_{n+1}=(12n^2+12n+4)u_n-32n^2u_{n-1}$ with $\lambda=(8,4)$, $k=2$,
$\xi=\tfrac12G$. By Proposition~\ref{prop:F4} its signed binomial transforms are
integral four-term rows on the same local system --- for $\nu=1,2,-1,-2$ they
have $\lambda=\{8-\nu,4-\nu,-\nu\}$, $k=2$ and scores
$-3.0986,-2.6931,-3.6094,-3.7918$, so the cusp move even improves the elementary
score by up to $+0.69$ nats \Verified{integrality to $n=459$}. But:

\begin{proposition}[\Proved{} and \Verified{$70$ digits}]\label{prop:D4}
The signed binomial transform $T$ of the companion is \emph{not} the companion of
the transformed row. Writing $L_0B=t$ for Zagier $\mathbf E$ and
$\tilde y(t)=(1+\nu t)^{-1}y(t/(1+\nu t))$ for the gauge, one has
$L'(Ty)(t)=(L_0y)(s)$ with $s=t/(1+\nu t)$, so $L'(TB)=t/(1+\nu t)\ne t$ and $TB$
fails the four-term recurrence at every $n\ge1$. The row's own companion $W$
\textup{(}$w_1=1$, $L'W=t$\textup{)} is $W=TB+Ty$ with
$(n+1)^2y_{n+1}=P_0y_n-Q_0y_{n-1}+\nu^n$, $y_0=y_1=0$, and the exact Casoratian
$(m+1)^2(a_mb_{m+1}-a_{m+1}b_m)=32^m$ gives
$\xi_\nu=\tfrac12G\cdot A(\nu/32)-B(\nu/32)$ with $A=\sum a_nt^n$,
$B=\sum b_nt^n$.
\end{proposition}

The four values $\xi_1=0.48630684566125056749\dots$,
$\xi_2=0.52037151700596857713\dots$, $\xi_{-1}=0.43391612361971143068\dots$,
$\xi_{-2}=0.41311716537104535942\dots$ are all unidentified against the battery
of \S\ref{sec:battery} \Verified{as stated}, and structurally they must be, being
values of the period and of its Eichler integral at an \emph{interior} point of
the disc; the control $\tfrac12G$ itself is recovered by the battery at $79$
digits. So the four-term world does contain integral rows built on the Catalan local
system --- the whole $\nu$-family --- but it is a four-point local system wearing
a fifth point: no new arithmetic, and a canonical Ap\'ery limit that is no longer
$\tfrac12G$. The contrast with \S\ref{sec:mixed}, where the fold is a simple root
of $\Rc$, the row is honestly four-term and the canonical limit \emph{is}
$\tfrac14G$, is what makes the mixed-class rows a genuinely new realisation
rather than a re-labelling.


% ===========================================================================
\section{The six-point glimpse: five-term rows}\label{sec:sixpoint}
% ===========================================================================

One term further, the five-term row
\[
(n+1)^2u_{n+1}=P(n)u_n-Q(n)u_{n-1}+R(n)u_{n-2}-T(n)u_{n-3},
\qquad u_0=1,\ u_{-1}=u_{-2}=u_{-3}=0,
\]
with $T=kn^2+ln+m$ has an operator with \emph{six} singular points and three
accessory parameters.

\begin{theorem}[five-term dictionary; \Proved, \Verified{as stated}]\label{thm:D5}
$L=\theta^2-tP(\theta)+t^2Q(\theta+1)-t^3R(\theta+2)+t^4T(\theta+3)
 =t^2\Rc D^2+t\Sc D+t\Vc$ with
\[
\begin{aligned}
\Rc&=1-at+dt^2-gt^3+kt^4,\\
\Sc&=1-(a+b)t+(3d+e)t^2-(5g+h)t^3+(7k+l)t^4,\\
\Vc&=-c+(d+e+f)t-(4g+2h+j)t^2+(9k+3l+m)t^3 .
\end{aligned}
\]
The exponents are $(0,0)$ at $0$, $(0,\rho_i)$ at the four roots of $\Rc$ with
$\rho_i=-T_c(t_i)/(t_i\Rc'(t_i))$,
$T_c=\Sc-t\Rc'=1-bt+(d+e)t^2-(2g+h)t^3+(3k+l)t^4$, and $3-s_1,3-s_2$ at $\infty$
with $s_i$ the roots of $T$; Fuchs reads $\sum_i\rho_i=s_1+s_2-2$. All four
exponents equal $\rho$ iff $b=(1-\rho)a$, $e=-2\rho d$, $h=-(1+3\rho)g$,
$l=-(2+4\rho)k$, and then $(j_1+j_2)/M=s_1+s_2=2+4\rho$ for
$T(n)=C(Mn-j_1)(Mn-j_2)$. The Casoratian on four solutions is
$W_{n+1}=-T(n)W_n/(n+1)^2$.
\end{theorem}

The free parameters are $a,c,d,f,g,j,C$: three accessory parameters $c=u_1$,
$f=Q(0)$, $j=R(0)$ --- exactly $6-3$ --- plus four position moduli modulo the
scaling. The scanner was validated against an exact brute force on a small box,
both finding the same $17$ integral rows, every one with a repeated
characteristic root \Verified{as stated}.

\subsection{The census and the disguise fraction}
Seven classes were scanned over $a\le24$, $|c|\le12$, $|d|\le90$, $|f|\le40$,
$|g|\le90$, $|j|\le60$, $|C|\le90$ in census mode, plus a deeper box on the
productive class. In the all-semistable class $(0;1,1,1)$, $T=C(n-1)^2$ (six
$I_n$ fibres), about $1700$ integral rows are found, $99.6\%$ with a repeated
characteristic root, leaving $\mathbf 7$ genuine six-point rows; in the gauge
partner $(-1;1,-1,-1)$ all $1608$ integral rows have a repeated root, leaving
$\mathbf 0$ \Verified{boxes as stated}. The pattern of \S\ref{sec:scanshape}
continues and intensifies: there the dominant degeneracy was the apparent
singularity ($98.2\%$ disguised three-term rows); here it is a repeated
characteristic root --- the singular points collide and the object is not a
six-point one at all --- at $99.6\%$--$100\%$. The gauge partner
$(-1;1,-1,-1)$ is the five-term analogue of $\mathbf S_{-1}$: empty, exactly as
Proposition~\ref{prop:F5} makes $\mathbf S_{-1}$ empty. The scan of $(0;1,1,1)$
was still running when it was stopped, so the list of seven is not claimed
complete \Open.

\begin{center}\footnotesize
\begin{tabular}{lllrl}
\toprule
$(a,c,d,f,g,j,C)$ & $\lambda$ & $k$ & score & $\xi$\\
\midrule
$(1,0,-14,-2,20,8,-8)$ & $2{\pm}2\sqrt2,\,-2,\,-1$ & $2$ & $-2.6931$ & $\tfrac{\sqrt2}3L(g,2)$\\
$(6,2,0,0,-6,-2,-1)$ & $3{\pm}2\sqrt2,\,-1,\,1$ & $2$ & $\mathbf{-2.0000}$ & $\tfrac{\sqrt2}3L(g,2)$ again\\
$(8,2,-52,-8,72,24,-16)$ & $12.5918,\,-2.3177,\,-2,\,-0.2741$ & $2$ & $-2.8406$ & $0.2336058844421865619710\dots$\\
$(4,1,-13,-2,9,3,-1)$ & $6.2959,\,-1.1588,\,-1,\,-0.1371$ & $2$ & $-2.1474$ & $0.4672117688843731239421\dots$\\
$(9,3,13,2,4,1,-1)$ & $7.2959,\,1,\,0.8629,\,-0.1588$ & $2$ & $\mathbf{-2.0000}$ & the same $0.46721176888\dots$\\
$(5,2,0,0,-20,-8,-16)$ & $4,\,-2,\,2,\,1$ & $2$ & $-2.6931$ & $\xi^*$ (\S\ref{sec:xistar})\\
$(5,2,17,3,31,12,18)$ & $1{\pm}2\sqrt2\,i,\,2,\,1$ & $2$ & $-3.0986$ & none (complex pair)\\
\bottomrule
\end{tabular}
\end{center}

$k=2$ for all seven --- the same free-integration phenomenon one term further ---
and \emph{every period occurs at least twice}. Three of the seven have
$|\lambda_2|=1$ exactly, hence score $-2.0000$.

\subsection{A weight-$3$ cusp-form $L$-value as an honest Ap\'ery limit}
The first row of the list,
\begin{equation}\label{eq:fiveterm}
(n+1)^2u_{n+1}=(n^2{+}n)u_n+(14n^2{+}2)u_{n-1}+(20n^2{-}20n{+}8)u_{n-2}
 +8(n{-}1)^2u_{n-3},
\end{equation}
$u_n=1,0,4,8,40,144,616,2544,10936,47200,\dots$, is integral \Verified{exact
rational arithmetic to $n=400$; $u_{400}$ has $271$ decimal digits}, primitive
($a=1$), and has all $u_n\ge0$ for $n\le400$. Its characteristic quartic is
$\lambda^4-\lambda^3-14\lambda^2-20\lambda-8=(\lambda+1)(\lambda+2)(\lambda^2-4\lambda-4)$
with $\Delta=2048\ne0$, so the six singular points
$0,\ -1,\ -\tfrac12,\ \tfrac{\sqrt2-1}2,\ -\tfrac{\sqrt2+1}2,\ \infty$ are
distinct; the exponents are $(0,0)$ at each of the five finite ones, so a
logarithm is forced and none is apparent, and $(2,2)$ at $\infty$. Sharp $k=2$
and $\text{score}=\log\tfrac12-2=-2.6931$. Its Ap\'ery limit is
\[
\xi=0.528749455595766833775958063859865871341\dots
=\frac{\sqrt2}{3}L(g,2)=\tfrac12\xi_{K3}(1)
=\re\,\xi_{K3}\Bigl(\tfrac{5\pm i\sqrt2}{27}\Bigr)
\]
\Verified{$110$ digits}, where $L(g,2)$ is the critical value of the weight-$3$
CM newform \texttt{32.3.d.a} of Theorem~\ref{thm:k3}.

This matters because the elliptic-$K3$ four-term row has a complex conjugate
dominant pair and therefore no archimedean Ap\'ery limit at all: there $L(g,2)$
appears only as a fold constant, obtained by analytic continuation. In the
six-point world the same $L$-value is an honest $\lim b_n/a_n$ of an integral
row, with a score $0.955$ nats better. The two local systems are not obviously
the same --- $\{5\pm i\sqrt2,1\}$ over $\Q(\sqrt{-2})$ against
$\{2\pm2\sqrt2,-2,-1\}$ over $\Q(\sqrt2)$, five points against six --- but
$\texttt{32.3.d.a}$ has CM by $\Q(\sqrt{-2})$ and level $32=2^5$, so both are
$2$-power-level objects and a correspondence is the obvious guess; identifying it
is open \Open, as is a five-term analogue of the $\Jm$-map test. Two further
facts: $t=\infty$ has exponents $(2,2)$ with a logarithm, so the operator has a
\emph{second MUM point}, whose characteristic quartic
$(\lambda-1)(\lambda-2)(\lambda^2-2\lambda-1)$ again has $|\lambda_2|=2$ and
whose slice contains no integral row \Verified{$|c|\le40$, $|f|\le60$,
$|j|\le80$}; and $|\lambda_2|=2$ exactly, so as a lattice engine the row's only
candidate alignment prime is $2$.

\subsection{Periods repeat across worlds}
Five independent numerical identities turned up, in two shapes. \emph{(A)} A
six-point Ap\'ery limit equals the real part of a five-point fold constant: the
row \eqref{eq:fiveterm} against the elliptic $K3$; the row $(6,2,0,0,-6,-2,-1)$
against the same $K3$ row, with score $-2.0000$ against its $-3.6479$; and the
row $(8,2,-52,-8,72,24,-16)$ against the mixed row $(\tfrac12,0;3,2,4)_{r=27}$,
$(23,8,-216,-42,-324)$ --- each to $110$ digits --- plus, inside the four-term
world, the cross-row identity of \S\ref{sec:folds} at $70$ digits.
\emph{(B)} One Ap\'ery limit shared by five rows in three classes and two
worlds: the constant $\xi^*$ of \S\ref{sec:xistar}, to $319$ digits.

\begin{conjecture}[\Conj]\label{conj:D6}
The real part of a complex fold constant of a five-singular-point four-term row
is, systematically, an honest archimedean Ap\'ery limit $\lim b_n/a_n$ of an
integral six-singular-point five-term row.
\end{conjecture}

The five instances are \Verified{$70$--$319$ digits}. The mechanism is
presumably a quadratic base change: separating a complex-conjugate pair of
singular points adds a puncture and replaces the fold constant by its
$\Gal(\C/\R)$-trace $\tfrac12(\xi_++\xi_-)$. No period found in this work occurs
only once, which suggests the census is seeing families rather than sporadic rows
--- but the identity is verified, not explained \Open.

% ===========================================================================
\section{Arithmetic: what does not change}\label{sec:arith}
% ===========================================================================

\subsection{No new irrationality candidate}
The \emph{score} of a row is $\log(1/|\lambda_2|)-k$ with $\lambda_2$ the second
largest characteristic root in modulus; a positive score is the archimedean part
of an irrationality proof. Two complete sweeps settle the question in the
four-term world.

Box W ($|\lambda_2|<1$; $a\le300$ for $18$ classes, $a\le400$ for $\mathbf S_0$)
contains $79$ rows with $|\lambda_2|<1$, all in the classes $(\pm1;3,\ldots)$ and
$(\pm1;4,\ldots)$, and \emph{every one has a vanishing Frobenius obstruction}: a
disguised three-term row.

Box PW ($|\lambda_2|<e^{-1}$, hence complete for a positive score whenever
$k\ge1$; all $27$ classes, swept to completion for $a\le800$, $|c|\le60$,
$|f|\le120$) contains \emph{exactly two integral rows in total}, both in the
class $(0;2,-1,3)$ with $\delta_\infty=2$:
\[
(144,32,96,-32,4):\ \lambda=143.331,\,0.3506,\,0.3184,\ k=1,\ \text{score}=+0.0480,
\]
\[
(225,50,-120,40,4):\ \lambda=225.532,\,-0.2662\pm0.0092i,\ k=1,\ \text{score}=+0.3229,
\]
and \emph{both have an exactly vanishing Frobenius obstruction at $t=\infty$}
(the value is $0$, not merely small); the exponents there are
$(\tfrac12,\tfrac52)$, so the monodromy is $-I$, an $I_0^*$ fibre invisible to
the projective local system, which therefore has only four singular points. They
are the four-term avatars of the classical Legendre--Pad\'e family, and $k=1$
identifies them as Pad\'e constructions.

\begin{corollary}[\Verified{boxes W and PW as stated}]\label{cor:F6}
Over the whole scanned region --- every Kodaira-admissible equal-exponent class,
both windows --- \emph{no genuine five-singular-point four-term row has a
positive score}. The only positive scores that occur belong to disguised
three-term rows with $k=1$, exactly as in the rigid case, where every positive
score outside Ap\'ery's and Beukers' rows has $k=1$ and is a classical
Legendre--Pad\'e row. The four-term world produces no new irrationality
candidate.
\end{corollary}

Theorem~\ref{thm:D1} adds the structural half: a decaying row needs an
irreducible characteristic cubic, so it cannot live in an unequal-exponent class,
and by Theorem~\ref{thm:D2} a real decayer needs a totally real one. Neither
condition is met anywhere in the scanned region.

\subsection{The $2$-adic content of the new Catalan rows}\label{sec:2adic}
The six rows of \S\ref{sec:mixed} are \emph{engines}, not decayers, so what
matters for the two-row lattice machinery is their $p$-adic arithmetic. Write
$\kappa_p=-\lim v_p(\operatorname{den}b_n)/n$,
$w_p=\lim v_p(\mathrm{Cas}_n)/n$ with $\mathrm{Cas}_n=a_nb_{n+1}-a_{n+1}b_n$, and
$\sigma_p=w_p+2\kappa_p$ the slope of $v_p(b_n/a_n-b_{n-1}/a_{n-1})$. Everything
below is exact (rational or $p$-adic), archimedean at $1500$ digits, measured to
$n=800$ for $p=2$ and $n=400$ for all $p$, with $p$ swept to $71$
\Verified{as stated}.

\begin{center}\footnotesize
\begin{tabular}{clrlrrll}
\toprule
row & $\xi_\infty$ & $g$ & $\kappa_p$ & $w_2$ & $\sigma_2$ & slope primes $p\le71$ & $\xi_2$\\
\midrule
1 & $\tfrac14G$ & $-128$ & $\kappa_2=-\tfrac12$, else $0$ & $\tfrac92$ & $\tfrac72$ & $\{2\}$ & $\tfrac14\zeta_2(2)$\\
2 & $\tfrac12G-\tfrac3{16}\zeta(2)$ & $8$ & $0$ & $2$ & $2$ & $\{2\}$ & $\tfrac12\zeta_2(2)$\\
3 & $\tfrac38\zeta(2)-\tfrac12G$ & $32$ & $0$ & $\tfrac72$ & $\tfrac72$ & $\{2\}$ & $-\tfrac12\zeta_2(2)$\\
4 & $\tfrac1{32}(2\zeta(2)+15L_{-3})$ & $32$ & $0$ & $3$ & $3$ & $\{2\}$ & $\tfrac38L_2(2,\chi_{12})$\\
5 & $\tfrac1{16}(15L_{-3}-6\zeta(2))$ & $16$ & $0$ & $3$ & $3$ & $\{2\}$ & $\tfrac34L_2(2,\chi_{12})$\\
6 & $\tfrac18(2\zeta(2)-3L_{-3})$ & $-1$ & $0$ & $0$ & $0$ & $\varnothing$ & none\\
Zagier $\mathbf E$ & $\tfrac12G$ & $32$ & $0$ & $5$ & $5$ & $\{2\}$ & $\tfrac12\zeta_2(2)$\\
\bottomrule
\end{tabular}
\end{center}

The $2$-adic limits are exact to $1172$--$1200$ $2$-adic digits, and the transfer
$G\mapsto\zeta_2(2)$, $\zeta(2)\mapsto0$, $L(2,\chi_{-3})\mapsto\tfrac45L_2(2,\chi_{12})$
is linear on the nose: rows 4 and 5 give $\tfrac{15}{32}\cdot\tfrac45=\tfrac38$
and $\tfrac{15}{16}\cdot\tfrac45=\tfrac34$, exactly as observed.

Three things are genuinely new in \emph{form} \Verified{as stated}.
\emph{Half-integral slopes}: rows 1 and 3 have $\sigma_2=\tfrac72$ and row 1 has
$w_2=\tfrac92$, whereas every row in the earlier census has an integer
$\sigma_p=v_p(c)$; a half-integral Casoratian slope means the two $2$-adic Newton
slopes of the pair $(a_n,b_n)$ are conjugate over a \emph{ramified} quadratic
extension --- here $\Q_2(\sqrt2)$, so the irrationality of the singular points is
visible $2$-adically. \emph{Decoupling}: $\sigma_2=\tfrac72,2,\tfrac72,3,3,0$
against $v_2(g)=7,3,5,5,4,0$, so the three-term law $\sigma_p=v_p(c)$ fails for
every one of the six and the slope must be measured; the $3\times3$ Casoratian
still obeys $W_{n+1}=R(n)W_n/(n+1)^2$ exactly, but the $2\times2$ one has no
closed form. \emph{A negative $\kappa_2$}: row 1 is an integral row whose
numerators are divisible by a growing power of $2$, $v_2(a_n)=\tfrac n2+O(1)$, so
$\kappa_2=-\tfrac12$ and $\sigma_2=w_2+2\kappa_2=\tfrac92-1=\tfrac72$ on the
nose --- the invariant $\sigma=w+2\kappa$ confirmed in a regime it was never
tested in.

But the \emph{content} is old. The slope set is $\{2\}$ for rows 1--5 and empty
for row 6 (checked to $p=71$), exactly as for Zagier $\mathbf E$, and the
$2$-adic values are the ledger's own two constants with new rational scalars: row
2's $\xi_2=\tfrac12\zeta_2(2)$ is literally $\mathbf E$'s value with $\sigma_2=2$
instead of $5$, and rows 4 and 5 sit in Zagier $\mathbf F$'s $2$-adic class but
lack $\mathbf F$'s second slope prime. The engine yield
$\sum_p\sigma_p\log p-\log|\lambda_2|$ is
$0.3466,\ 0.6931,\ 1.7329,\ 0.0693,\ 2.0101,\ 0$
for rows 1--6, against $2.0794$ for Zagier $\mathbf E$ and $\log9=2.1972$ for
$\mathbf F$: \emph{every one of the six is dominated by the engine it would
replace}.

\subsection{The Catalan ledger}
The obstruction analysis \cite{CATOBS} lists, as the evasion (E1) of the
geometric wall, ``a host with an irrational second singularity'': three-term rows
cannot have one, their characteristic roots being rational hence integral, so
$|t_2|\le\tfrac14$ at every placement and the conformal size is too small for the
Hodge depth; rows with more singular points can. The six rows of
\S\ref{sec:mixed} realise (E1) explicitly --- irrational conjugate singularities,
five punctures, periods in $\Q G+\Q\zeta(2)$ at $139$ digits, full classical
period matrices, best score $+0.693$ nats over Zagier $\mathbf E$ and equal to
the best Catalan geometry previously known --- and the ledger does not move: the
$2$-adic slope set is still $\{2\}$ and the limits still lie in the $\zeta_2(2)$
and $L_2(2,\chi_{12})$ classes, so the second wall (2.3) of \cite{CATOBS}, that
the $p$-adic places cannot compensate because $G\in\Q$ is false at $p=2$, is
untouched. The Catalan class's deficit is a missing $p=3$ partner, and rows
1--3 supply $\{2\}$ again with a \emph{smaller} $2$-adic resource than
$\mathbf E$; meanwhile Theorem~\ref{thm:D1} forbids a real decayer in any
unequal-exponent class. In the language of \cite{CATOBS}: (E1) has been
realised, and (2.3) has been re-decorated, not evaded. Whether the
conditional-function machinery applies to rows whose fold is an order-$2$
orbifold point rather than a cusp is not settled here \Open.

% ===========================================================================
\begin{remark}
The modular Catalan hosts themselves are not four-term rows: on the level-$16$ and $\Gamma_0(12)$ hosts the fold is a \emph{double} root of the reversed characteristic polynomial, and the minimal recurrences have $6$, $7$ and $6$ terms; the cusp move destroys their Ap\'ery limits (Proposition D4 of \texttt{FOUR\_TERM\_DEEP.md}, $\xi_\nu=\tfrac12G\,A(\nu/32)-B(\nu/32)$). So the four-term world and the modular Catalan world intersect only through the mixed-exponent rows of \S4.
\end{remark}

\section{Open problems}\label{sec:open}
% ===========================================================================

\begin{enumerate}[label=\textup{\arabic*.},leftmargin=2.2em]
\item \emph{Integrality of the sporadic rows: Zagier's conjecture one step out.}
For three-term rows integrality is explained by modular parametrisation, and the
unconditional classification is exactly the statement that every integral row has
one \cite{CLASS,Zagier2009}. Here $29$ of the $30$ rows of
Table~\ref{tab:thirty} fail the $\Jm$-map test with $\deg\le18$, and no mechanism
is known for their integrality; even for the one modular row, the elliptic $K3$
\eqref{eq:k3row}, integrality is experimental, and we do not know whether the
index-$16$ group attaining the Riemann--Hurwitz bound is congruence \Open.
\item \emph{Identify $\xi^*$} (\S\ref{sec:xistar}): one transcendental shared by
five rows in three classes and two worlds, to $319$ digits, unidentified against
$145$ classical constants. The natural next step is the fold constants of the
$+\tfrac12$ gauge and of the $(0,\mp\tfrac12;2,\cdot)$ classes, not computed
here; if the pattern of \S\ref{sec:matrix} holds they should lie in
$\Q\,\Gamma(\tfrac14)^4/\pi+\Q\cdot(\text{something})$, which would pin $\xi^*$
\Open.
\item \emph{Classify five-point elliptic surfaces with integral rows.} The
configurations with five singular fibres form one-parameter families, so the
locus they cut out in the $(\Ic_1,\Ic_2)$-plane is a curve for each fibre type;
computing those curves would give a genuine analogue of the cross-ratio
finiteness of the rigid case, and would \emph{predict} the $K3$ row instead of
discovering it. The enumeration of the fifteen signatures of \S\ref{sec:RH} into
conjugacy classes of genus-zero subgroups of $\PSL_2(\Z)$ --- the analogue of the
$28$ groups in the four-point case --- is also not carried out \Open.
\item \emph{Conjecture~\ref{conj:D6}}: prove the base-change mechanism behind the
five identities, and in particular explain why a six-point row over $\Q(\sqrt2)$
has the period of a five-point elliptic $K3$ over $\Q(\sqrt{-2})$, and why the
second such row does it with a score $1.65$ nats better than the $K3$ row's own
\Open.
\item \emph{Are the six mixed rows a family?} The class $(-\tfrac12,0;1,0,0)$ is
complete over the box of \S\ref{sec:xistar} and contains exactly nine primitive
rows, at $r=\pm1,\pm2,\pm3,4,8,8$, so within that box they are sporadic; whether
they are the small members of a one-parameter family at larger $|r|$ or $|a|$ is
open, and would change the picture, the three Catalan periods being three points
of the two-dimensional space $\Q G+\Q\zeta(2)$ that a family would sweep out
\Open.
\item \emph{Loose ends:} a closed form for $v_2(\mathrm{Cas}_n)$ of the six rows
(measured as $w_2n+O(1)$ with the $O(1)$ in $[-3,3]$, exact at multiples of
$100$) and a proof of $v_2(a_n^{(1)})=\tfrac n2+v_2(a_n^{\mathbf E})$, observed
at eleven checkpoints; a $\Jm$-map test adapted to the mixed and five-term
classes, none of which has been tested for modularity; the $21$ equal-exponent
Ap\'ery limits and the $24$ new fold periods, which are probably weight-two
\emph{multiple} polylogarithms at arguments in $\Q(\lambda_i)$, of which only the
classical and depth-one parts were tested; the unequal-exponent classes with
$\rho_1\ne\rho_2\ne\rho_3$, which need their own parametrisation; and
completeness --- box PW is complete for the positive-score question and box W for
the decayer question in the stated classes and ranges, but no deep box is, and
the cost law of \S\ref{sec:boxes} makes a literal $100\times$ extension of box G
about $4\times10^6$ core-seconds \Open.
\end{enumerate}

\subsection*{Reproducibility}\label{sec:repro}
All statements above are reproduced from the project documents in
\texttt{consolidation/}: \texttt{FOUR\_TERM\_SCAN.md} (the dictionary, the
classes, the census, the $30$ rows, the $K3$ row's discovery,
Corollary~\ref{cor:F6}); \texttt{FOUR\_TERM\_DEEP.md} (Theorems~\ref{thm:D3},
\ref{thm:D1}, \ref{thm:D2} and \ref{thm:D5}, the mixed classes, the six Catalan
rows and their period matrices, the gauge partners, $\xi^*$, the $2$-adic
census, the five-term probe, Conjecture~\ref{conj:D6});
\texttt{K3\_ROW\_PERIOD.md} for Theorem~\ref{thm:k3}; and
\texttt{CATALAN\_OBSTRUCTION.md} for the ledger discussion of
\S\ref{sec:arith}. Scripts, job files and logs --- the
scanners and their brute-force validators, the exact re-verification, the
$\Jm$-map test, the fold-constant machinery, the period batteries and the
$\kappa_p/w_p/\sigma_p$ census --- are under the directories
\texttt{lattice/four\_term/}, \texttt{lattice/four\_term\_deep/} and
\texttt{lattice/k3\_period/}; the output files carry every hit with its verdict
and invariants.

\enlargethispage{3\baselineskip}
\begin{thebibliography}{99}
\scriptsize
\setlength{\itemsep}{0pt plus 0.5pt}

\bibitem{AlmkvistZudilin2006}
G.~Almkvist and W.~Zudilin,
\emph{Differential equations, mirror maps and zeta values},
in: Mirror Symmetry V, AMS/IP Stud. Adv. Math. \textbf{38} (2006), 481--515.

\bibitem{Beauville1982}
A.~Beauville,
\emph{Les familles stables de courbes elliptiques sur $\PP^1$ admettant quatre
fibres singuli\`eres},
C. R. Acad. Sci. Paris S\'er. I Math. \textbf{294} (1982), 657--660.

\bibitem{Beukers2002}
F.~Beukers,
\emph{On Dwork's accessory parameter problem},
Math. Z. \textbf{241} (2002), 425--444.

\bibitem{CLASS}
Claude (Fable) and River,
\emph{Integral second-order Ap\'ery-like recurrences and rational elliptic
surfaces}, companion paper (\texttt{paper/classification}), 2026.

\bibitem{K3PAPER}
Claude (Fable) and River,
\emph{An Ap\'ery-like sequence attached to a singular $K3$ surface},
companion note (\texttt{paper/k3\_row}), 2026.

\bibitem{CATOBS}
Claude (Fable) and River,
\emph{Catalan's constant and the arithmetic holonomy method: an obstruction},
companion paper (\texttt{paper/catalan\_cdt}), 2026; project document
\texttt{consolidation/CATALAN\_OBSTRUCTION.md}.

\bibitem{Cooper2012}
S.~Cooper,
\emph{Sporadic sequences, modular forms and new series for $1/\pi$},
Ramanujan J. \textbf{29} (2012), 163--183.

\bibitem{Herfurtner1991}
S.~Herfurtner,
\emph{Elliptic surfaces with four singular fibres},
Math. Ann. \textbf{291} (1991), 319--342.

\bibitem{Kodaira1963}
K.~Kodaira,
\emph{On compact analytic surfaces II, III},
Ann. of Math. \textbf{77} (1963), 563--626; \textbf{78} (1963), 1--40.

\bibitem{Livne}
R.~Livn\'e,
\emph{Motivic orthogonal two-dimensional representations of
$\mathrm{Gal}(\overline{\Q}/\Q)$},
Israel J. Math. \textbf{92} (1995), 149--156.

\bibitem{LMFDB}
The LMFDB Collaboration,
\emph{The $L$-functions and modular forms database},
\texttt{https://www.lmfdb.org}, 2026; newform \texttt{32.3.d.a}.

\bibitem{MirandaBook}
R.~Miranda,
\emph{The basic theory of elliptic surfaces},
ETS Editrice, Pisa, 1989.

\bibitem{MovasatiReiter}
H.~Movasati and S.~Reiter,
\emph{Heun equations coming from geometry},
Bull. Braz. Math. Soc. (N.S.) \textbf{43} (2012), 423--442.

\bibitem{Persson1990}
U.~Persson,
\emph{Configurations of Kodaira fibers on rational elliptic surfaces},
Math. Z. \textbf{205} (1990), 1--47.

\bibitem{K3P}
\texttt{FOUR\_TERM\_SCAN.md}, \texttt{FOUR\_TERM\_DEEP.md},
\texttt{K3\_ROW\_PERIOD.md} (project documents, \texttt{consolidation/}, 2026).

\bibitem{Zagier2009}
D.~Zagier, \emph{Integral solutions of Ap\'ery-like recurrence equations},
CRM Proc. Lecture Notes \textbf{47} (2009), 349--366.

\end{thebibliography}

\end{document}
